%% LyX 2.3.7 created this file.  For more info, see http://www.lyx.org/.
%% Do not edit unless you really know what you are doing.
\documentclass[english]{extarticle}
\PassOptionsToPackage{natbib=true}{biblatex}
\usepackage{lmodern}
\usepackage[T1]{fontenc}
\usepackage[latin9]{inputenc}
\usepackage[a4paper]{geometry}
\geometry{verbose,tmargin=1in,bmargin=1in,lmargin=1in,rmargin=1in}
\setlength{\parskip}{\medskipamount}
\setlength{\parindent}{0pt}
\usepackage{babel}
\usepackage{float}
\usepackage{amsmath}
\usepackage{graphicx}
\usepackage{setspace}
\onehalfspacing
\usepackage[unicode=true]
 {hyperref}

\makeatletter

%%%%%%%%%%%%%%%%%%%%%%%%%%%%%% LyX specific LaTeX commands.
%% Because html converters don't know tabularnewline
\providecommand{\tabularnewline}{\\}

%%%%%%%%%%%%%%%%%%%%%%%%%%%%%% User specified LaTeX commands.
\usepackage{graphicx}
\date{}

\@ifundefined{showcaptionsetup}{}{%
 \PassOptionsToPackage{caption=false}{subfig}}
\usepackage{subfig}
\makeatother

\usepackage[style=apa]{biblatex}
\addbibresource{biblio_syit.bib}
\begin{document}
\title{Stress\nobreakdash-Testing Survey\nobreakdash-to\nobreakdash-Survey
Imputation: Understanding When Poverty Predictions Can Fail\thanks{The authors acknowledge financial support from the World Bank. We
thank Chris Elbers, Roy van der Weide, Sergio Olivieri, Hai-Anh Dang,
and Nobuo Yoshida for comments. We also thank participants of the
S2S workshop held at the World Bank. Additionally, we thank Luis-Felipe
Lopez Calva, Gabriela Inchauste, and Maria Eugenia Genoni for providing
support and space to pursue this work. Thanks are also due to Tara
Vishwanath for her insights and guidance on earlier background work.
Special thanks to Ankriti Singh wit DECPR for ensuring replicability
of the paper's results. Any error or omission is the authors\textquoteright{}
responsibility alone.}}
\author{Paul Corral, Andres Ham, Peter Lanjouw, Leonardo Lucchetti, and Henry
Stemmler\thanks{Paul Corral, Leonardo Lucchetti, and Henry Stemmler are part of the
World Bank. Andres Ham is of Universidad de los Andes, and Peter Lanjouw
is part of Vrije Universiteit Amsterdam. Corresponding author: Paul
Andres Corral Rodas, \protect\href{http://pcorralrodas@worldbank.org}{pcorralrodas@worldbank.org}.}}

\maketitle
\vfill{}

\textbf{Key words:} Poverty, Inequality, Poverty imputation, missing
data 

\textbf{JEL classification: }I32, C53

\section{Introduction}

The measurement and monitoring of poverty are central to assessing
global development progress. For the World Bank, whose mission is
to eradicate extreme poverty on a livable planet, the ability to track
poverty reduction is fundamental for measuring institutional effectiveness
and guiding policy decisions. Yet, a persistent challenge hampers
this crucial task: the limited availability of recent, high-quality
household survey data that includes information that allows for comparable
measures of expenditure or income across time. This challenge is particularly
acute in countries where poverty is most concentrated \citep{DangJolliffeCarletto,worldbank2024poverty}. 

The issue affects many countries, and hinders monetary poverty monitoring
due to the lack of data in some of the world's most populous and poverty
afflicted economies. For example, India and Nigeria account for a
substantial share of the world's poor, meaning that they heavily influence
global poverty trends. In India, official consumption survey data
was unavailable between 2011 and 2022.\footnote{The 75th round of the NSO Survey of Consumption Expenditure was collected
in 2017-18 but was not released until after 2024: \href{https://www.mospi.gov.in/unit-level-data-report-nss-75th-round-july-2017-june-2018-schedule-250social-consumption-health}{https://www.mospi.gov.in/unit-level-data-report-nss-75th-round-july-2017-june-2018-schedule-250social-consumption-health}.} In Nigeria comparable survey data for poverty monitoring was absent
between 2009 and 2018, and as of 2025 the country has no new publicly
available consumption survey data. These data gaps not only affect
our understanding of poverty at the country level but compromise our
ability to produce reliable global poverty estimates and assess progress
toward poverty reduction goals.

Traditional poverty measurement relies on household surveys that collect
detailed consumption or income data. These surveys represent substantial
investments in both financial and human resources. They require extensive
preparation, careful implementation, and can place considerable burden
on responding households, who must either maintain detailed consumption
diaries or participate in comprehensive recall interviews. The method
of data collection itself can introduce significant biases into poverty
estimates which may compromise comparability of estimates over time.\footnote{For instance, consumption diaries -- while theoretically more accurate
-- can lead to respondent fatigue and underreporting over time. Recall
modules, on the other hand, may suffer from memory bias, with longer
recall periods typically resulting in lower reported consumption.
Research has shown that simply changing the recall period or the number
of consumption items can lead to substantially different poverty estimates
within the same population (\citealt{beegle2012methods}).} These methodological challenges add another layer of complexity to
the already demanding task of poverty measurement.

The challenges of collecting traditional living standards surveys
for poverty measurement become particularly acute during crises when
standard survey operations are disrupted or impossible to conduct.
Armed conflicts, natural disasters, health emergencies, and other
humanitarian crises often prevent face-to-face household interviews
precisely when poverty monitoring becomes most crucial. The COVID-19
pandemic provided a stark illustration of this challenge, as household
survey efforts were halted globally at a time when understanding welfare
impacts was most critical. Similar data collection constraints arise
in conflict zones, where security concerns prevent enumerator access
to households, or during natural disasters that displace populations
and disrupt statistical operations. During such crises, policymakers
and international organizations must resort to alternative methods
for estimating poverty. 

Survey-to-survey (S2S) imputation has emerged as a methodological
response to these data limitations. The method is often used as a
way to produce poverty numbers when traditional welfare data collection
is unavailable due to costs, complexity, conflict, or other reasons.
S2S, often combined with rapid phone surveys or other alternative
data collection methods, has emerged as a key tool for maintaining
poverty monitoring during these challenging periods where traditional
surveys cannot be collected. 

S2S methods build upon techniques originally developed for small area
estimation (SAE) in poverty mapping, pioneered by \citet{jesko1998}
and further refined by \citet{elbers2003micro}. These imputation
methods enable poverty estimation using surveys that lack direct welfare
measures by:
\begin{enumerate}
\item Developing a predictive model of household welfare using a survey
with consumption or income data (the \textquotedbl source\textquotedbl{}
survey) 
\item Applying this model to a different survey containing similar household
characteristics but lacking direct welfare (the ``target'' survey)
\item Generating poverty estimates based on the predicted welfare distribution
\end{enumerate}
While S2S and small area estimation share methodological foundations,
they serve distinct purposes. Unit-level small area estimation of
poverty, as developed by \citet{elbers2003micro} and advanced by
\citet{molina2010small}, typically applies nationally estimated models
(from a household survey) to population census data to generate precise
poverty estimates for small geographic areas where survey estimates
are often too imprecise for reliable estimates. In contrast, S2S focuses
on applying the model to a separate survey to predict welfare, regardless
of geographic disaggregation. One might view traditional poverty mapping
as a special case of S2S where the target datasets are the different
areas of the census and the primary goal is obtaining geographically
disaggregated estimates.

Although a potentially an attractive solution, S2S approaches come
with their own limitations and potential biases that must be carefully
considered. This paper first examines the limitations and potential
pitfalls of survey-to-survey imputation through rigorous analysis
of its fundamental assumptions. This is a technical paper that relies
on simulated data to conduct experiments. In these experiments, the
focus is not on traditional concerns such as variable selection or
model specification, but rather on the method's core limitations,
particularly for measuring poverty across extended time periods or
different populations. The analysis relies on simulated data to isolate
and demonstrate specific mechanisms that can generate biased estimates.
This approach allows for clear illustration of three critical findings: 
\begin{enumerate}
\item Standard sampling bias-correction techniques, such as re-weighting
to match population means, may be insufficient when source and target
surveys differ fundamentally. 
\item The method's tendency to replicate the welfare distribution of the
source survey can make it unreliable for measuring changes in inequality
and poverty.\footnote{The model captures the joint distribution between $y$ and $x$ in
the source data. The parameters corresponding to the source data are
then applied to the target data, where the only difference is the
distribution of covariates, $x$. Thus, unless all the change is captured
by changes in the covariates the method will likely give biased estimates.} Thus, better communication behind the limitations of the methods
is crucial. These insights are particularly relevant given the increasing
reliance on survey-to-survey imputation to fill data gaps in poverty
monitoring. 
\item Third, omitted variable bias will likely affect poverty predictions,
particularly when imputing across time periods that include significant
economic shocks or structural changes. This last finding has important
implications for applications that use \textquotedbl fast-moving\textquotedbl{}
variables to capture welfare changes,\footnote{\citet{SwiftMethod} refer to expenditure related covariates as fast
moving. These include dichotomous variables indicating if a household
reported purchasing a product or not.} as these variables may introduce bias if they are correlated with
unobserved factors that affect welfare.
\end{enumerate}
Through systematic examination of these limitations using simplified
examples, the paper provides practitioners with a framework for understanding
when and why S2S methods can fail. The results here are of relevance
since survey-to-survey imputation has become more prevalent in recent
times to overcome data scarcity.


\section{The effectiveness of S2S across time: A summary of the evidence \label{sec:Survey-to-Survey-Imputation}}

Estimating poverty relies on household survey data that collects household
consumption or income. Collecting household consumption and income
data is costly and complex. For consumption, modules typically span
multiple domains---food, housing, health, education, and more---each
requiring detailed recall and careful categorization. These surveys
are usually fielded over a full year to represent seasonal consumption
patterns. As \citet{deaton1998consumption} observe, these modules
are among the most demanding components of household surveys, often
encompassing hundreds of items and multiple recall periods. This not
only increases respondent burden but also drives up implementation
costs. Ensuring accuracy adds another layer of complexity: respondents
may forget or misreport expenditures, and cultural and regional differences
often require localized adaptations in survey design. Moreover, as
\citet{beegle2012methods} show, even small changes in how surveys
are administered can yield significant differences in reported poverty
and consumption levels.

Because of the complexity and cost, few countries collect comparable
consumption or income data frequently. Long-term, harmonized consumption
trends are the exception, not the norm. Yet, policymakers, donors,
and development practitioners continually need up-to-date poverty
statistics. To fill this gap, survey-to-survey imputation has become
an increasingly used method for estimating poverty when consumption
or income data are unavailable.

S2S builds upon the idea proposed by \citet{elbers2003micro,jesko1998}
where a linear regression model is fit on a survey where consumption
or income are available to determine the joint distribution of consumption
or income, and a given set of household characteristics -- this survey
is typically referred to as the source or training data. The model
parameters are then applied to a separate dataset (the target data)
where the same set of household characteristics can be found with
the aim of replicating the welfare distribution (consumption or income)
via imputation. With a simulated welfare distribution in hand, poverty
estimates and other welfare indicators can be derived from the target
data. The methods can be applied to contemporaneous surveys or to
surveys corresponding to a different point in time. 

S2S is useful when applied to contemporaneous surveys to obtain statistics
across the welfare distribution for indicators that are not available
in the living standards survey. For example, the Demographic and Health
Surveys (DHS) provide a litany of health related indicators that are
not available in most living standards surveys that collect welfare
for poverty estimation. This is where S2S comes into play. A model
trained on the living standards survey and applied to the DHS can
help us inform about the difference in stunting rates between poor
and non-poor households or determine how it varies across the welfare
distribution. 

The method is often treated as a prediction exercise and unlike much
of econometrics, the parameters of the linear regression are not intended
to capture only the direct effect of the characteristics on welfare
(\citealt{elbers2003micro}, p356). Nevertheless, the original intent
of \citet{elbers2003micro} was to predict on to the same population. 

When imputing to a different year, the method rests on the assumption
that model parameters remain stable---that is, the relationship between
household characteristics and welfare has not changed. In other words,
the method assumes that model parameters remain stationary over time
-- meaning that any observed changes in poverty are solely attributable
to changes in the model's covariates, rather than shifts in unobservable
factors or changing returns to these covariates (Dang, Lanjouw and
Serajuddin \citeyear{dang2017}). As \citet{christiaensen2012small}
note, such assumptions may be especially problematic in rapidly growing
economies, where structural economic changes can alter the relationship
between poverty and its predictors. It can also be problematic when
a strong shock has occurred between the survey periods -- for example,
a global pandemic or changes in a country's fragility. \citet{mathiassen2009model}
emphasizes that the validity of imputations rests on a strong assumption
of parameter stability and explicitly warns against extrapolating
too far across time or applying the method when significant changes
in welfare dynamics or survey design are likely.

Researchers and practitioners have proposed various approaches to
address the limitations of S2S in predicting poverty over time. \citet{stifel2007tracking}
advocate for including time-varying variables such as rainfall and
prices to capture temporal changes. \citet{SwiftMethod} suggest incorporating
variables that track economic conditions more directly, like indicators
on whether the household consumed a given good. However, as \citet{SwiftMethod}
emphasize, without updated training data to re-estimate model parameters,
the risk of missing significant economic changes remains substantial,
even with these additional covariates. \citet{dang2025poverty} further
note that evidence on the utility of item-specific consumption indicators
(e.g., expenditures or quantities for particular goods) remains inconclusive,
and instead advocate for parsimonious models that are easier to implement
in surveys and in challenging field conditions. 

The effectiveness of S2S imputation across time has been tested in
a range of contexts, with mixed results. For example, Dang, Lanjouw
and Serajuddin (\citeyear{dang2017}) evaluate S2S using Jordan\textquoteright s
Household Expenditure and Income Survey and the Employment and Unemployment
Survey; \citet{stifel2007tracking} conduct similar tests for Kenya;
and Dang et al. (\citeyear{dang2025poverty}) apply the method across
multiple countries---including Ethiopia, Malawi, Nigeria, Tanzania,
and Viet Nam---reporting estimates generally within acceptable error
margins. While some studies report estimates within acceptable margins
of error (\citealt{dang2025poverty}, \citealt{https://doi.org/10.1111/roiw.12007}),
others find significant discrepancies between predicted and observed
poverty rates (\citealt{newhouse2014survey}, \citealt{https://doi.org/10.1111/roiw.12007}).
A useful example comes from Malawi, where \citet{mathiassen2021predicting}
document their efforts to produce comparable poverty estimates between
2004 and 2014. While their method performed well within individual
surveys and with careful seasonal adjustment, comparability across
survey rounds proved fragile. Shifts in survey design and implementation
such as shorter questionnaires, periods of data collection, changing
field protocols, and others led to unstable predictor relationships
and inconsistent poverty trends over time. Even in this relatively
well-structured case, the authors conclude that S2S imputation should
only be used when comparability of predictors and survey instruments
can be assured (\citealt{mathiassen2021predicting}).

\citet{christiaensen2012small} undertake an empirical validation
of survey-to-survey imputation methods over time. The authors perform
survey-to-survey imputation in scenarios where there is comparable
expenditure data which provides a ``true'' estimate of poverty.
The authors validate their approach with data for Viet Nam and for
China using rural household panel data. In Viet Nam, the authors obtain
a model using the 1992/93 data and predict poverty using the 1997/98
data. They note that the method works relatively well and depending
on the covariates used, differences between predicted and observed
poverty rates were on average 3.4 percentage points during a period
where poverty fell by 23.2 percentage points. For the Chinese regions
where the method was tested, the authors also find that the methods
work relatively well. However, depending on the model used, differences
between predicted and observed rates were considerable. This highlights
the method's conundrum: different models may generate different findings,
and at least some models may get it right. However, ex-ante, it is
very difficult to determine which model will yield the most accurate
estimates or even the right estimates. So far, no foolproof methods
or diagnostics exist to reliably identify the ideal model, raising
fundamental concerns about the robustness of imputed poverty estimates.

Applications of survey-to-survey (S2S) imputation have also made their
way into global poverty monitoring---for example, until recently,
every poverty number for India after 2011 was based on an imputation.
However, evidence from real-world applications, particularly when
time gaps are large or covariate distributions change substantially,
reveals mixed findings and highlights the need for caution. These
findings underscore the importance of understanding the context and
data very well when applying S2S across different periods and economic
environments. To better appreciate how such estimates are produced
and the assumptions they rest on, the next section outlines the foundations
of survey-to-survey imputation, tracing its origins in the broader
field of missing data methodologies and clarifying its application
in the context of poverty measurement.

\section{The basics behind survey-to-survey imputation (S2S)\label{sec:The-basics-behind}}

Imputation is a method for filling in missing data. According to \citet{van2018flexible},
the first instance of a statistical method to replace a missing value
dates to 1930,\footnote{\citet{allan1930method}.} and the first widespread
use of the term ``imputation'' comes from \citet{madow1983incomplete}.\footnote{As noted by \citet{van2018flexible}.}
The method was originally proposed to fill in missing observations
and considered the nature of the missing data (\citealt{dempster1977maximum}).
The goal of multiple imputation is not to create a single prediction,
but rather multiple imputations to reflect the uncertainty around
the actual value (\citealt{van2018flexible}). By generating plausible
values for missing data, multiple imputation allows for valid estimation
and the construction of appropriate confidence intervals \citep{van2018flexible}.
For example, in the case of a regression on agricultural yields, the
variable capturing plot sizes may have missing values and these must
be imputed with the aim of not losing information as well as obtaining
a valid estimate of the coefficient for the relationship between land
and yields. Using multiple imputation in this example, as opposed
to just the predicted land size, is expected to reduce the rate of
false positives.\footnote{Type I errors, which occur when a null hypothesis is incorrectly rejected.} 

The case of a variable that is entirely missing in the dataset was
not considered by the original multiple imputation literature. The
academic background for predicting an entirely missing variable in
a given dataset is more aligned to the small area estimation literature
(see \citealt{rao2005small} and \citealt{rao2015small}). In model-based
small area estimation auxiliary information from censuses, registers
or other large surveys are leveraged to produce statistics with sufficient
precision where survey data alone is not enough, for example small
geographical areas. Perhaps the first instance of survey-to-survey
imputation, as is applied in this paper, comes from the work of \citet{jesko1998}
which is more aligned to small area estimation and noted by the authors.
The key difference to small area estimation up to that point was that
\citet{jesko1998} predicted the variable of interest, consumption,
at the household level and from that they obtained aggregate statistics
based on the prediction of consumption. While earlier examples exist
of combining data from different sources to predict household-level
variables (e.g., Arellano and Meghir \citeyear{arellano1992female}),
the innovation in \citet{jesko1998} lies in applying models estimated
from survey data to census data to replicate the full consumption
distribution. This approach allows for the estimation of poverty and
other welfare indicators as if welfare data were available in the
census itself. After refinements by \citet{elbers2003micro}, the
work became the basis of what later was referred to as poverty mapping
in the World Bank. The key difference between survey-to-survey imputation
and small area estimation is that the latter aims to replicate the
welfare distribution for each specific area or group of interest,
rather than only for the entire population. One of the first applications
of the methods from \citet{elbers2003micro} to predict national level
poverty can be seen in the work of \citet{simler2004estimating} who
use the methods to track changes in poverty in Mozambique.

Survey-to-survey (S2S) imputation for poverty measurement relies on
the assumption that a population's welfare distribution can be captured
by a linear model which is estimated on the $source$ data. For simplicity,
the subscript $target$ is used to indicate when parameters are applied
to the target data. Following the presentation by \citet{molina2010small},
but ignoring the nested error component,\footnote{This is later considered for robustness checks.}
the population's assumed data generating process (DGP) for transformed
welfare $\ln y_{i}$ is:\footnote{Transformed data is often used as the dependent variable to ensure
the model's assumptions hold. For simplicity, throughout this document
the assumed transformation is the natural logarithm, although many
others are possible. Meeting the model's statistical assumptions is
crucial for reliable estimation. \citet{corral2022guidelines} demonstrate
that violations of these assumptions, particularly the normality of
residuals, can lead to biased poverty estimates. Their analysis shows
that appropriate transformation of the dependent variable (household
welfare) can significantly reduce such bias. When the standard logarithmic
transformation proves insufficient, alternative transformations may
be necessary to better approximate normality in the model's error
term.} 

\begin{equation}
\ln\,y_{i}=x_{i}\beta+e_{i};\:i=1,\ldots,N\label{eq:welfare model}
\end{equation}

\[
e\stackrel{i.i.d.}{\sim}N(0,\sigma_{e}^{2})
\]

where $x_{i}$ is a vector of independent variables common to the
source survey and the target survey, and $e_{i}$ is a random disturbance
term that is assumed to be distributed i.i.d. and $N(0,\sigma_{e}^{2})$.
The $\beta$ as well as the $\sigma_{e}^{2}$ parameters are estimated
using the source survey or the training data since this is the only
data where the welfare vector of interest ($\ln\,y_{i}$) is available.
These estimated parameters are then applied to the target data to
predict $\ln\,y_{i}$, and from that poverty or other welfare related
indicators.

Because normally distributed errors are assumed in Eq. \ref{eq:welfare model},
although this can be relaxed as later illustrated, for any given household
$i$ in the target survey, the probability of being poor for household
$i$ is entirely dependent on its expected welfare, $x_{i}\beta,$
and the distribution of its idiosyncratic term, $e_{i}$, which is
assumed to follow $e_{i}\sim N(0,\sigma_{e}^{2})$ 

\begin{equation}
\Pr\left(poor_{i}=1\right)=\Phi\,\left(\frac{\ln z-x_{i}\beta}{\sqrt{\sigma_{e}^{2}}}\right)\label{eq:cum prob}
\end{equation}

where $\ln z$ is the natural log of the chosen poverty threshold,\footnote{Note that if the transformation of the dependent variable is not the
natural logarithm, then this is not valid. The poverty line must be
transformed in a similar manner as the dependent variable.} and $\Phi\left(.\right)$ is the standard normal CDF. Consequently,
what the method calculates is each household's probability of being
poor. The average probability of being poor across households corresponds
to the national poverty rate.\footnote{Note that because the only thing differing across households in Eq.
\ref{eq:cum prob} are the characteristics, $x_{i}$, a proxy means
test (PMT) approach that relies only on $x_{i}\beta$ will yield the
same household ranking as the survey-to-survey approach if the same
model is used for either. Under a PMT, the threshold is chosen at
a given percentile of $x\beta$ so that by construction it yields
a desired proportion of people eligible.} 

The implementation of welfare predictions to the target data can follow
different approaches. The most common method, grounded in the multiple
imputation literature, generates several welfare vectors to reflect
prediction uncertainty. This approach is aligned to what was implemented
in poverty mapping through the PovMap software (\citealt{zhao2006user})
and is similar to the methodology used in Stata's multiple imputation
commands (\texttt{mi regress}). Under this approach, each imputed
welfare vector is generated through a five-step process:\footnote{Taken from \texttt{mi regress} \citet{StataCorp2023}.}
\begin{enumerate}
\item Obtain model parameter estimates for $\beta$, and $\sigma_{e}^{2}$
via ordinary least squares fit on the training data (Eq. \ref{eq:welfare model}
), noted below as $\hat{\beta}$ and $\hat{\sigma}_{e}^{2}$.
\item The distribution of the error terms is drawn from its posterior Chi-square
distribution, where $n$ is the number of observations in the training
or source data which is sampled from $N$, and $K$ is the number
of covariates in the model, including the model intercept:
\[
\sigma_{e}^{2*}\sim\hat{\sigma}_{e}^{2}\frac{\left(n-K\right)}{\chi_{n-K}^{2}},
\]
\item With $\sigma_{e}^{2*}$ in hand it is possible to update the variance-covariance
matrix needed for drawing the $\beta^{*}$ parameters: 
\[
\beta^{*}\sim MVN\left(\hat{\beta},\sigma_{e}^{2*}\left(X'X\right)^{-1}\right),
\]
\item Then errors are drawn from:
\[
e_{i}^{*}\sim N\left(0,\sigma_{e}^{2*}\right)
\]
\item Each imputed transformed welfare for household $i$ in the $target$
data under imputation $m$ is given by: 
\begin{equation}
\ln y_{im}^{*}=x_{i_{target}}\beta_{m}^{*}+e_{im}^{*}\label{eq:MI}
\end{equation}
\item Repeat steps 2 through 5 to obtain $M$ sets of imputed values, $\ln y_{i1}^{*},\ln y_{i2}^{*},\ldots,\ln y_{iM}^{*}$
\end{enumerate}
Hence, for each imputed vector, a new $\sigma_{e}^{2*}$ is drawn
and used to draw $\beta^{*}$ and a household specific residual, $e_{i}^{*}$.
The multiple imputation simulation approach is not necessarily aligned
to reducing the prediction mean squared errors (MSE), but as noted
by \citet{van2018flexible} the purpose is to minimize the likelihood
of false positives, since normally the imputed vectors are frequently
used for regression analysis (see \citealp{rodas2021pull} for a discussion
related to SAE). Alternatively, step 4 may be skipped and the imputed
transformed welfare can be obtained from $\ln y_{im}^{*}\sim N\left(x_{i_{target}}\beta_{m}^{*},\sigma_{e}^{2*}\right).$

An alternative approach, derived from the small area estimation literature,
treats the source survey as the best representation of the true welfare
distribution. Under this approach, parameters estimated from the source
survey are applied directly to the target data through Monte Carlo
simulation. For straightforward indicators like poverty rates, practitioners
can simply apply the asymptotic formula for expected values (Eq. \ref{eq:cum prob}).
The steps are straight forward:
\begin{enumerate}
\item Fit Equation \ref{eq:welfare model} to the training data using ordinary
least squares to obtain estimates of the model parameters, $\hat{\beta}$
and $\hat{\sigma}_{e}^{2}$.
\item Each transformed welfare for household $i$ in the Monte Carlo simulation
is given by -- note how the coefficients remain constant across Monte
Carlo simulations:\footnote{Monte Carlo simulations are not always necessary and for poverty one
could immediately apply Eq.\ref{eq:cum prob}. However, for more complex
parameters Monte Carlo simulations allow for indicators that are difficult
to handle analytically.}
\begin{equation}
\ln y_{im}^{*}\sim N\left(x_{i_{target}}\hat{\beta},\hat{\sigma}_{e}^{2}\right)\label{eq:Fixed}
\end{equation}
\item Repeat step 2 to obtain $M$ sets of imputed values, $\ln y_{i1}^{*},\ln y_{i2}^{*},\ldots,\ln y_{iM}^{*}$
\end{enumerate}
Recent methodological advances have expanded model fitting options
to include machine learning (ML) techniques. However, since most ML
methods do not make explicit assumptions about error term distributions,
they cannot directly replicate the multiple imputation approach described
above. Nevertheless, bootstrap-based alternatives exist where draws
from the empirical residuals from the source survey are used to obtain
a welfare vector ($e_{i}^{*}=\hat{e}_{j},$where $j\sim\mathrm{Uniform}\{1,2,\ldots,n\}$).
These were implemented in both PovMap (\citealt{zhao2006user}) and
Stata's \texttt{sae }command by \citet{nguyen2018sae}. This bootstrapping
approach has proven particularly valuable for regularized regression
techniques, as demonstrated by \citet{lucchetti2024application} with
lasso regression,\footnote{\href{https://github.com/pcorralrodas/lassopmm}{https://github.com/pcorralrodas/lassopmm}.}
and may be extended to other ML methods such as gradient boosting,
random forests, or Bayesian additive regression trees. 

An additional method that can be paired with machine learning techniques
as well as OLS, is predictive mean matching (PMM). The approach compares
the predicted values from the target data to those from the training
data, specifically by minimizing the absolute difference between linear
predictors $|x_{i_{target}}\hat{\beta}-x_{i}\hat{\beta}|$. Once a
specified number of nearest neighbors is identified, one of the actual
observed welfare values is randomly drawn \citep{van2018flexible}.
PMM is notably less sensitive to model misspecification and is robust
to back transformation (ibid). For instance, models estimated on $\log(y)$
typically yield similar results to those based on $\exp(y)$ (\citealt{van2018flexible}),
although differences in explained variance can still affect poverty
estimates. The number of donors (i.e. neighbors) is determined by
the practitioner. While setting this number to one can result in repeated
values---especially when model fit is weak---\citet{van2018flexible}
recommends using five neighbors. He also emphasizes that this choice
should be adaptive, depending on the sample size and the model\textquoteright s
explanatory power, as measured by its $R^{2}$.


\subsection{Imputing to a contemporaneous survey\label{subsec:sametime}}

Imputation of poverty to a contemporaneous survey can be a useful
tool. For example, assume a country collects a living standard survey
annually and has also collected a Demographic and Health Survey (DHS).
If one would like to determine how nutritional indicators among children
vary across the monetary welfare distribution, survey to survey imputation
may be of use to impute a monetary welfare aggregate to the DHS to
determine what share of children in the poorest quantiles are stunted.
Nevertheless, practitioners could also impute the indicator of interest
to the living standards survey. For example, with the goal of determining
energy poverty in Bulgaria, \citet{rude2024statistical} impute energy
spending from the household budget survey (HBS) to the European Survey
of Income and Living Conditions to determine profiles of the energy
poor. These settings resemble the \textquotedblleft no single dataset
contains all variables\textquotedblright{} problem discussed by \citet{rubin1986statistical},
where imputations may be performed in either direction to enable inference. 

A fundamental limitation, however, is that when key variables such
as stunting and welfare are never jointly observed, their conditional
association, given the common covariates used for matching, is not
identified from the observed data. As a result, identifying stunting-by-welfare
gradients requires an additional, inherently untestable restriction
(most commonly conditional independence) or a sensitivity analysis
that varies the assumed association between stunting and welfare (\citealp{Raessler2004DataFusion}).

The successful application of survey-to-survey imputation depends
on meeting several preconditions, first outlined by \citet{jesko1998}
and later refined by \citet{elbers2003micro}. These preconditions
ensure that the fundamental assumption -- both surveys represent
the same underlying population -- holds. Two contemporaneous surveys
which represent the same population are needed for the approach to
work well. In this instance, the assumption of stable parameters does
not affect the method's estimates.

First, covariate comparability: the predictors used to impute welfare
must be available in both surveys and defined and measured in a consistent
way. This extends beyond the presence of similarly named variables
to encompass harmonized concepts, response categories, and field protocols;
differences in definitions (for example, how household membership
is determined) can alter both the distribution of covariates and their
relationship to welfare.

Second, distributional alignment: the marginal distributions of key
predictors should be similar across the source and target surveys,
not only in means but also in dispersion. Large discrepancies may
indicate that the target survey places substantial weight on regions
of the covariate space with limited support in the source survey,
undermining the model\textquoteright s ability to replicate the welfare
distribution.

Third, structural coherence: the joint relationships among predictors---and,
by implication, the mapping from predictors to welfare---should be
broadly consistent across surveys. In a contemporaneous setting, parameter
instability over time is not the primary concern; rather, the key
risk is that survey design differences or measurement inconsistency
of covariates shift underlying associations in ways that make the
welfare model non-portable.

\textbf{Model Assumptions }

The reliability of welfare predictions in the same period depends
on meeting the model's statistical assumptions. Conditional on comparability
and alignment across surveys, the quality of welfare imputations hinges
on the adequacy of the statistical model used to generate predicted
welfare ($x'\hat{\beta}$) and simulated disturbances ($e^{*}$). 

Three assumptions are particularly salient. First, many implementations
rely on approximately normal residuals; when this is violated, transformations
of the welfare aggregate can improve approximation to normality (\citealp{corral2022guidelines}).
If transformations remain insufficient, an alternative is to draw
disturbances from the empirical residual distribution, though this
requires additional structure---most notably that residuals are appropriately
centered and, in common implementations, approximately symmetric.\footnote{Although OLS residuals are mean-zero in sample, symmetry matters when
residuals are used as a stand-in for the unobserved component in simulation.
If the residual distribution is skewed, drawing errors unconditionally
from the empirical residuals can systematically alter the estimated
probabilities. If large positive values are more likely to be sampled
than large negative values, or vice versa, it is consequential for
poverty prediction. } 

Second, the model often assumes homoscedastic errors; this can be
relaxed by explicitly modeling heteroskedasticity, as in \citet{elbers2002micro}
and \citet{harvey1976estimating}. Third, the mapping from covariates
to welfare is frequently treated as linear, but can be made more flexible
through richer parametric specifications (e.g., polynomial terms and
interactions) or through nonparametric and machine-learning estimators
such as random forests or gradient boosting.

\subsection{Imputing across time}

Survey-to-survey (S2S) imputation is promoted as a valuable tool for
filling gaps in welfare data, but its use over time requires greater
caution than for contemporaneous applications as discussed in section
\ref{sec:Survey-to-Survey-Imputation}. The method assumes that relationships
estimated in a given period can be mapped to a different point in
time. The core assumption---that changes in welfare reflect only
shifts in observable characteristics, with stable relationships between
covariates and welfare---becomes increasingly fragile as the time
between surveys grows. Structural economic change, evolving labor
markets, and shifting household behaviors can all violate this assumption,
leading to imputation errors that distort actual welfare dynamics.
Moreover, S2S methods tend to reproduce the welfare distribution of
the source survey: coefficients capture the joint distribution of
$y$ and $x$ from the source period, and errors are drawn from that
same model in the source data. As a result, households with identical
characteristics in the source and target periods are assigned identical
poverty probabilities, regardless of contextual changes. Because the
method assumes stable parameters, it cannot reflect how the poverty
risk of households with similar observed characteristics may change
over time, which limits its ability to recover true changes in poverty
and inequality.

Additionally, omitted variables can be particularly consequential
when imputing across time. In contemporaneous settings, these concerns
are sometimes viewed as second-order for prediction because, in the
estimation sample, OLS residuals are mean-zero and the intercept absorbs
the average influence of unobservables, so fitted values can perform
well in reproducing mean welfare (see \citealt{a4f1365d-ff19-34b5-b1b5-15ec4c1e5f15,elbers2003micro,dang2025poverty}).\footnote{\citeauthor{dang2025poverty} (2025 p9) states that ``worries about
endogeneity of $x_{j}$ pose far less important, if any, concerns
in our context.'' An endogenous variable is frequently defined as
an explanatory variable that may be correlated with the error term
(Wooldridge, \citeyear{wooldridge2009introductory} p88). Omitted
variables are related as these are correlated to the error term and
a covariate.} This rationale is less compelling in cross-time applications, where
the distribution of unobserved determinants may shift between the
source and target periods due to shocks, policy changes, or structural
transformation. In that case, source-period parameters can embed period-specific
correlations between observables and unobservables that are not portable,
implying biased welfare and poverty predictions even when covariates
are measured consistently.\footnote{See Annex \ref{subsec:Omitted-Variables} for an exposition.}

These issues are instances of a broader concern: parameter instability.
Cross-time S2S applies both $\beta$ and the residual structure (e.g.,
$\sigma_{e}^{2}$ or a modeled heteroskedastic variance) estimated
historically to a different period. When returns evolve---as emphasized
in decompositions that allow coefficients to change over time (\citealt{bourguignon2008beyond})---imputations
may primarily reflect compositional change rather than the true evolution
of the welfare distribution. This can be seen from $\beta=(X'X)^{-1}X'y$:
substantial shifts in the distribution of $X$ change both $X'X$
and $X'y$, with no guarantee that the implied mapping remains valid
in the target period. In practice, cross-time reliability can be undermined
by changes in determinants and returns, sampling frames,\footnote{Changes in sampling frame are not necessarily problematic, as long
as the data is still representative of the population of interest.} policy regimes, macroeconomic shocks, and survey design or implementation.

Finally, these limitations are not resolved by selecting models solely
to optimize in-sample fit (e.g., via stepwise selection or regularization).
Portability depends on the stability of the conditional relationship
being transported; if regressors proxy for unobserved determinants
whose correlation structure changes over time, relationships that
are predictive within-period can become unreliable out-of-period.
Given the growing reliance on S2S to monitor poverty in the absence
of frequent welfare aggregates, cross-time applications should be
interpreted cautiously and, where feasible, triangulated against independent
signals of welfare dynamics---such as national accounts aggregates,
labor-market indicators, price statistics, or relevant administrative
series.

\section{Model-Based Simulations}

This section builds on the concepts introduced in Section \ref{sec:The-basics-behind}
by examining how violations of key survey-to-survey (S2S) assumptions
translate into bias and instability in imputed poverty estimates.
The core analysis relies on simulated data, which provides a controlled
environment in which individual features of the data-generating process
(DGP) can be perturbed in isolation. This design allows the analysis
to attribute changes in performance to specific assumption failures---something
that is difficult to do with observational data, where multiple sources
of non-comparability and model misspecification tend to move together.

To strengthen the external relevance of the simulation findings, two
robustness extensions are introduced.\footnote{We thank anonymous reviewers for these recommendations.}
First, an alternative DGP is considered in which the \textquotedblleft linear\textquotedblright{}
component of the welfare model (i.e. $x\beta$) is far away from normally
distributed, providing a stricter test of approaches that rely on
distributional assumptions. Second, the analysis uses empirical covariate
structure from Mexico following \citet{corral2021map}: the linear
predictor is estimated on the Mexican intracensal data and then combined
with simulated error processes that vary in form---(i) nested normal
errors consistent with clustered disturbance structures, (ii) heteroskedastic
errors, and (iii) non-normal errors generated from a Student\textquoteright s
$t$ distribution. Together, these exercises are intended to evaluate
whether the paper\textquoteright s recommendations remain robust when
residual behavior departs from the baseline assumptions in ways that
are empirically plausible in household survey settings.

Model-based simulations thus serve a dual purpose in what follows:
they provide a transparent framework for stress-testing S2S imputation
under targeted assumption violations, and they support more defensible
guidance by demonstrating which findings persist across alternative
DGPs and empirically grounded error structures.

\subsection{Creating populations\label{subsec:Creating-populations}}

\subsubsection{DGP relying on real world data\label{subsec:DGP-Mex}}

To complement the stylized simulations, the paper also evaluates robustness
using covariate structure estimated from real world data. Specifically,
the publicly available Mexican Intercensus Survey of 2015, carried
out by the Mexican National Institute of Statistics and Geography
(INEGI), is used. The survey consists of 5.9 million households and
is representative across 2,457 municipalities of the country and localities
with a population larger than 50,000 people. It also includes an income
measure of welfare.

The analysis uses the processed data described in \citet{corral2021map}
and \citet{corral2025poverty}. There, the original intercensus data
were modified to randomly remove 90 percent of households that did
not report income, additionally all municipalities with less than
500 households sampled were removed. The final census consists of
3.9 million households across 1,865 municipalities. From this frame,
501 survey samples of roughly 23,540 households each are drawn, with
10 households per PSU, reflecting design features common in multi-stage
household surveys \citep{grosh1996manual}. Each sample includes all
32 states and between 951 and 1,020 municipalities. Further details
on the sampling strategy are provided in \citet{corral2021map}.

To construct the experimental data used here, a 10 percent simple
random sample across all PSUs in the census is taken. This subsample
is used to estimate a linear model of log income using the full set
of eligible household- and municipality-level covariates, following
\citet{corral2021map}. The resulting model's linear fit ($X\hat{\beta}$)
is kept fixed, but different disturbances are added to generate $\ln y$:
\begin{enumerate}
\item Nested normal error model: independent normally distributed disturbances
at the PSU, municipality, and household levels, calibrated to the
estimated variance components in the data to ensure the 
\item Non normal errors: error terms are drawn from a Student's $t$ distribution
with 8 degrees of freedom, scaled to ensure an SD of roughly 0.8,
and a model $R^{2}$ of roughly 0.45.
\item Multiplicative heteroskedasticity: A covariate, $x_{het}\sim N(0.5,0.5^{2})$,
is simulated. Then conditional on $x_{het}$ the error terms are drawn:
$e_{het}\sim N\left(0,\sigma^{2}(x_{het})\right),$where $\sigma=\frac{\exp\left((1/6)x_{het}\right)}{1.368}$
to ensure that the model's RMSE is roughly 0.8 and an $R^{2}$ or
roughly 0.45.
\end{enumerate}
The observed welfare variable is not used directly in these experiments.
The purpose of the Mexican exercise is not to reproduce Mexico-specific
welfare results, but to anchor the simulations in an empirically realistic
covariate structure while retaining control over the error-generating
process. This design isolates the performance consequences of the
disturbances noted above in a setting with realistic covariate distributions
and sampling features, avoiding conclusions that would otherwise be
driven by country- and period-specific features of the observed welfare
aggregate.

\subsubsection{DGP 1 - approximately normal $x\beta$\label{subsec:DGP-1--}}

We create 500 synthetic populations of 20,000 households where the
welfare of the population is generated with the following data generating
process (DGP): 

\begin{equation}
\ln y_{i}=3+0.1x_{1_{i}}+0.5x_{2_{i}}-0.25x_{3_{i}}+0.2x_{4_{i}}-0.15x_{5_{i}}+e_{i}\label{eq:dgp_model}
\end{equation}

where $e_{i}\sim N(0,0.5^{2})$
\begin{enumerate}
\item $x_{1}$ is a discrete variable, simulated as the rounded integer
value of the maximum between 1 and a random Poisson variable with
mean $\lambda=4$. It is rounded to create a variable similar to the
often used household size.
\item $x_{2}$ is a binary variable, taking value 1 when a random uniform
number between 0 and 1 is less than 0.2
\item $x_{3}$ is a binary variable, taking value 1 when a random uniform
number between 0 and 1 is less than 0.5 as long as $x_{2}=1$, otherwise
it is equal to 0
\item $x_{4}\sim N(2.5,2^{2})$
\item $x_{5}$ is a variable drawn from a Student\textquoteright s t distribution
with 5 degrees of freedom and scaled by 0.25
\end{enumerate}
This DGP yields a Gini coefficient of approximately 0.38, and the
covariates explain about 47 percent of the variation in $\ln y_{i}$. 

As an additional robustness check, Annex \ref{subsec:DGP_cluster}
considers an alternative DGP which follows an assumed double nested
error model similar to that of small area estimation based on \citet{ea7f7e4f-7381-39b0-abdd-8da14c150408},
while maintaining the covariate-generation process above. This extension
is used to illustrate that, when the true DGP includes clustered disturbances,
ignoring the nesting primarily affects inference and uncertainty rather
than systematically shifting bias and MSE in mean poverty predictions.

\subsubsection{DGP 2 - non-normally distributed $x\beta$\label{subsec:DGP-2--}}

A second DGP modifies the construction of $x_{4}$ and $x_{5}$ to
generate an $X\beta$ that departs from normality, while keeping the
generation of $x_{1},\:x_{2}$ and $x_{3}$ unchanged. Additionally,
$e_{i}\sim N(0,0.3^{2})$ is also adjusted to ensure the covariates
explain roughly 43 percent of the variation of $\ln y_{i}$ while
yielding a Gini coefficient of roughly 0.23. Specifically, $x_{4}$
and $x_{5}$ are simulated as follows:
\begin{enumerate}
\item $\tilde{x}$ is a categorical variable, that is simulated as the rounded
integer value of the maximum between 1 and a random Poisson variable
with mean $\lambda=5$
\item $x_{4}$ indicates low values of $\tilde{x}$ and is equal to 1 if
$\tilde{x}<3$.
\item $x_{5}$ indicates high values of $\tilde{x}$ and is equal to 1 if
$\tilde{x}\geq8$.
\end{enumerate}
Note that $\tilde{x}$ is not explicitly part of the model. Following
the previous section, a double nested error model structure is also
considered for this DGP.

\subsection{Evaluation metrics \label{subsec:Evaluation-metrics}}

The main performance metrics are bias and mean squared error (MSE).
Under the DGPs described in Section \ref{subsec:DGP-1--} and \ref{subsec:DGP-2--},
and unless otherwise noted, $B=500$ populations $\left(b=1,\ldots,B\right)$
are generated. For each population $b$, the imputation model is estimated
on a training sample and applied to a separate target sample to produce
an estimate of the indicator of interest, $\tau$ (head count monetary
poverty).

Empirical bias and MSE are computed across replications as:

\[
Bias\left(\tau\right)=\frac{1}{B}\sum_{b=1}^{B}\left(\hat{\tau}_{b}-\tau_{b}\right)
\]

\[
MSE\left(\tau\right)=\frac{1}{B}\sum_{b=1}^{B}\left(\hat{\tau}_{b}-\tau_{b}\right)^{2}
\]

where $\tau_{b}$ denotes the true population value under population
$b$, and $\hat{\tau}_{b}$ is the corresponding estimate obtained
by fitting the model on the training sample and predicting welfare
in the target sample. Note that under this process, following the
assumed DGP, $\tau$ is a random variable (\citealp{molina2010small}
p376).

For the Mexican data, the validation design differs slightly. The
population parameter $\tau$ is treated as fixed (observed in the
census-based file), and repeated survey samples are drawn from this
same population. A model is first estimated in a randomly selected
training sample and then applied to $S=500$ target samples drawn
from the population. In the small area estimation literature, this
approach is commonly referred to as design-based validation (\citealp{tzavidis2018start}).
Empirical design-based bias and MSE are computed as:

\[
Bias\left(\tau\right)=\frac{1}{S}\sum_{s=1}^{S}\left(\hat{\tau}_{s}-\tau\right)
\]

\[
MSE\left(\tau\right)=\frac{1}{S}\sum_{s=1}^{S}\left(\hat{\tau}_{s}-\tau\right)^{2}
\]

where $\tau$ denotes the true population value as observed in the
census and $\hat{\tau}_{s}$ is the corresponding estimate obtained
by applying the training-sample model to target sample $s$, for $s=1,\ldots,S$. 

\subsection{How to Impute?}

The typical imputation approach undertaken for S2S work done to predict
poverty follows the literature on multiple imputation (MI). A similar
method was followed in the original software implementation of small
area estimation proposed by \citet{elbers2003micro}, PovMap (\citealt{zhao2006user}).
Under the MI approach, the parameters estimated on the training data
are not applied directly to the target data, instead the parameters
to be applied to the target data are drawn from their posterior distributions
as illustrated in section \ref{sec:The-basics-behind}. For small
area estimation, the imputation approach from \citet{elbers2003micro}
has been updated to follow the approach from \citet{molina2010small}.\footnote{See a detailed discussion in \citet{rodas2021pull,corral2021map}.}
Under the method of the latter, the parameters estimated using the
training data are applied directly to the data and noise of the predicted
indicator is estimated via a parametric bootstrap which is aligned
to the model's assumptions \citep{gonzalez-manteiga2008bootstrap}.

To compare different imputation methods, a simple random sample (SRS)
comprising 20 percent of the data generated in Section \ref{subsec:DGP-1--}
is used as both source and target data. In the results presented in
this section that use the simulated data, the same dataset is used
for both model development and prediction. This approach eliminates
potential additional noise introduced by sampling variability, allowing
for a more controlled assessment of imputation performance that is
entirely driven by the model fit and its assumptions. The process
is applied across 500 generated populations allowing us to assess
the method's empirical bias. Additionally, where noted, results using
the Mexican data are presented. Under these simulations, models are
estimated in a given sample and then applied across 500 samples.

\subsubsection{Why simulate the errors?}

Survey-to-survey (S2S) imputation of poverty is not simply about predicting
the dependent variable. Relying solely on the predicted value of welfare
to derive poverty rates often yields poor results. This section uses
the synthetic population introduced in Section \ref{subsec:Creating-populations}
to illustrate this point.

We fit a linear model on a simple random sample (SRS) comprising 20
percent of the population. This model is then used to generate two
sets of predicted values for the full data:
\begin{enumerate}
\item The linear prediction, $\hat{y_{i}}=x_{i}\hat{\beta}$
\item The linear prediction plus a randomly drawn error: $y_{i}^{*}=x_{i}\hat{\beta}+e_{i}^{*}$,
where $e_{i}^{*}$ is drawn from $N(0,\hat{\sigma}^{2})$. 
\end{enumerate}
Figure \ref{fig:why_errors} shows that relying solely on the linear
prediction, $x_{i}\hat{\beta},$ fails to replicate the true distribution
of welfare. For example, if the true poverty line is 2.39---corresponding
to a 20 percent poverty rate in the actual data---then the linear
prediction alone would imply a poverty rate of only 10 percent. This
gap persists despite the fact that the linear prediction has a lower
mean squared error (MSE = 0.25 for both DGP) than the version with
stochastic error (MSE = 0.5 for both DGP).\footnote{The MSE here is measured as the mean squared difference in predictions
across all households. For the former the difference is $\left(y_{i}-x_{i}\hat{\beta}\right)$
and for the latter the difference is $\left(y_{i}-x_{i}\hat{\beta}-e_{i}^{*}\right).$ } This suggests that minimizing MSE of predicted welfare is not sufficient
for reproducing the distributional properties needed to estimate poverty
accurately. Moreover, the Spearman rank correlation is higher for
the linear prediction ($x\hat{\beta}$ vs. $y$) than for the version
with added error ($x\hat{\beta}+e^{*}$ vs. $y$). In other words,
correctly ranking households is also not enough to replicate the full
welfare distribution or derive accurate poverty statistics. 

\begin{figure}[H]
\caption{Replicating the actual distribution, with and without simulated errors
- DGP \ref{subsec:DGP-1--} (left), DGP \ref{subsec:DGP-2--} (right)
\label{fig:why_errors}}

\begin{centering}
\includegraphics[scale=0.62]{C:/Users/WB378870/GitHub/S2S_RTM/5.figures/1.Normal/why_errors}\includegraphics[scale=0.62]{C:/Users/WB378870/GitHub/S2S_RTM/5.figures/2.Nonnormal/why_errors}
\par\end{centering}
{\footnotesize{}Note: Data are generated as described in \ref{subsec:Creating-populations}.
True values correspond to those generated following Eq. \ref{eq:dgp_model},
XB are the predicted linear fit on to the full data, and XB+e is the
predicted linear fit plus a random error following the estimated distribution
from the model.}{\footnotesize\par}
\end{figure}

This underscores a central point: accurate estimation of poverty through
S2S imputation requires more than good point predictions or rankings.
It requires replicating the underlying distribution of welfare. 

\subsubsection{Under a normal error DGP\label{subsec:normal error dgp}}

The baseline imputation methods assume that welfare is linearly related
to a set of characteristics and that errors follow a two-fold double
nested normal error. The main results are based on a normal double
nested error distribution as illustrated in section \ref{subsec:DGP-Mex}.

Varying the number of imputations (20 versus 100) implemented using
Stata's \texttt{mi regress} has virtually no effect on bias or MSE
(Figure \ref{fig:bias_eb_mi}). However, the MI approach yields estimates
that are both more biased and noisier than approaches that treat model
parameters as fixed---namely the OLS and nested-error implementations
in Equation \ref{eq:Fixed}. These fixed-parameter predictions follow
the small area estimation approach of \citet{molina2010small}, where
$\beta$ and variance model parameters are kept fixed.

Ignoring the clustered nature of the data has a minimal impact on
the estimates' empirical bias and MSE. Point estimates that ignore
the nested structure (OLS) perform similarly to, and in this case
better than, estimates that explicitly model location effects (nested).
In small area estimation, \citet{ea7f7e4f-7381-39b0-abdd-8da14c150408}
suggest that modeling cluster effects at a level misaligned with the
level of reporting can yield biased noise estimates. \citet{corral2021map}
reinforce this point using model-based simulations but show that the
bias of point estimates remains unaffected when the model ignores
additional nested errors. In our case, because the results reported
here are at the national level, omitting random location effects does
not materially reduce the validity or precision of the poverty estimates.
These results are similar under the alternative DGPs described in
section \ref{subsec:DGP-1--} and \ref{subsec:DGP-2--}.\footnote{See annex \ref{subsec:DGP_cluster}.}

\begin{figure}[H]
\caption{Bias (left) and MSE (right) in FGT0 under a two-fold double nested
normal error DGP \label{fig:bias_eb_mi}}

\begin{centering}
\includegraphics[scale=0.62]{C:/Users/WB378870/GitHub/S2S_RTM/5.figures/3.Mexican/bias_normal_mex}\includegraphics[scale=0.62]{C:/Users/WB378870/GitHub/S2S_RTM/5.figures/3.Mexican/mse_normal_mex}
\par\end{centering}
{\footnotesize{}Note: Data are generated using Mexican data as described
in \ref{subsec:Creating-populations}. Bias is assessed at various
poverty lines across the welfare distribution. Specifically, these
lines correspond to percentiles that are multiples of 5. }\textbf{\footnotesize{}OLS}{\footnotesize{}:
uses the method described in Eq. \ref{eq:Fixed} to generate predictions,
ignoring the nested error structure. }\textbf{\footnotesize{}Nested:
}{\footnotesize{}similar to OLS but considers double nested error
structure of the DGP. }\textbf{\footnotesize{}MI X}{\footnotesize{}:
implements X imputations following Eq. \ref{eq:MI}. The model is
trained on one sample taken from the Mexican data, and applied to
500 different samples taken from the data. Poverty thresholds are
determined using the census.}{\footnotesize\par}
\end{figure}


\subsubsection{Under non-normally distributed errors}

It is essential to ensure the imputation method chosen is the one
best aligned to the data at hand. Data transformations, to ensure
the model assumptions are met can help (\citealt{rojas2020data,corral2021map}),
but there are instances where transformations are of little use (Figure
\ref{fig:sim_dgp_nonnormal_error}). If the residuals do not follow
a normal distribution, an alternative is to draw from the empirical
residuals. Stata's \texttt{mi regress} includes a bootstrap option
that estimates posterior parameters from bootstrap samples, addressing
concerns about asymptotic normality when parameter assumptions are
questionable (\citealp{StataCorp2023} \texttt{mi impute regress}
p2.). Similarly, the \texttt{hetmireg} command,\footnote{\texttt{hetmireg} (Corral - mimeo) command supports heteroskedasticity
following the alpha model described in \citet{elbers2002micro}. Corral
(mimeo) - https://github.com/pcorralrodas/hetmireg.} inspired from PovMap methods (\citealt{zhao2006user}), generates
bootstrap samples with errors drawn from the empirical distribution,
though it omits the area random effect required for area-level estimates
often necessary for small area estimation. 

Recent advancements have introduced machine learning (ML) techniques
for estimation, but since most ML methods lack explicit assumptions
about error term distributions, they cannot be directly applied. Instead,
bootstrap-based alternatives, implemented in PovMap (\citealt{zhao2006user})
and Stata\textquoteright s \texttt{sae} command \citep{nguyen2018sae},
have proven effective, particularly for regularized regression techniques
like lasso regression as shown by \citet{lucchetti2024application}.
These methods could be extended to other ML approaches, including
gradient boosting, random forests, and Bayesian additive regression
trees. \citet{verme2024predicting} offers a useful primer on the
performance of random forests and other ML approaches in poverty prediction.

\begin{figure}[H]
\caption{Bias (left) and MSE (right) in FGT0 of different methods under non-normal
errors for the Mexican data \label{fig:bias_eb_mi-1}}

\begin{centering}
\includegraphics[scale=0.62]{C:/Users/WB378870/GitHub/S2S_RTM/5.figures/3.Mexican/bias_nonnormal8_mex}\includegraphics[scale=0.62]{C:/Users/WB378870/GitHub/S2S_RTM/5.figures/3.Mexican/mse_nonnormal8_mex}
\par\end{centering}
{\footnotesize{}Note: Data are generated as described in \ref{subsec:DGP-Mex};
errors are simulated from a Student's t-distribution with 8 degrees
of freedom and scaled for a SD of 0.8. Bias and MSE are assessed as
noted in Section\ref{subsec:Evaluation-metrics} for the Mexican data
across poverty lines. Specifically, these lines correspond to percentiles
that are multiples of 5 taken from the full population. }\textbf{\footnotesize{}OLS
ignoring non-normal}{\footnotesize{}: uses the method described in
Eq. \ref{eq:Fixed} to generate predictions. }\textbf{\footnotesize{}OLS
BS 20}{\footnotesize{}: Estimates obtained using }\texttt{\footnotesize{}hetmireg,
}{\footnotesize{}OLS fit across 20 bootstrap samples of training sample
with errors drawn from the empirical distribution across bootstrap
samples. }\textbf{\footnotesize{}OLS Empirical}{\footnotesize{}: uses
the method described in Eq. \ref{eq:Fixed}, but errors are drawn
from the empirical distribution of the training sample. }\textbf{\footnotesize{}MI
BS}{\footnotesize{} }\textbf{\footnotesize{}20}{\footnotesize{}: Uses
Stata's bootstrap option for multiple imputation with 20 imputations.
}\textbf{\footnotesize{}Lasso Empirical: }{\footnotesize{}applies
parameters from lasso model where $\lambda$ is selected using 10
fold cross validation, and errors are drawn from the empirical distribution
of the training sample. }\textbf{\footnotesize{}RF empirical}{\footnotesize{}:
applies parameter estimates from a random forest model to the target
data, and errors are drawn from the empirical distribution of the
training sample.}{\footnotesize\par}
\end{figure}

Results from a simulation where the DGP for the Mexican data is adjusted
and errors are simulated from a Student's t-distribution with 8 degrees
of freedom and scaled for a SD of 0.8 are presented in Figure \ref{fig:bias_eb_mi-1}.
All methods perform well with bias that is smaller than 1 percentage
point across all poverty thresholds. With this considered, ignoring
the fact that errors are not normally distributed leads to estimates
that are more biased and noisier (OLS ignoring non-normal). Similarly,
Stata's \texttt{mi regress }with the bootstrap option does not yield
considerable improvement under this simulation (MI BS 20).\footnote{Stata's documentation for the option is unexpectedly short and provides
few details on the method's implementation.} The approaches that seem to work better in this scenario are those
based on the \texttt{hetmireg} command that relies on bootstrapped
samples with errors drawn from the empirical distribution of residuals
under each bootstrap sample of the training data (OLS BS 20).\footnote{The OLS implementation relies on the \texttt{hetmireg} command by
Corral (mimeo) - https://github.com/pcorralrodas/hetmireg.} Additionally, a more straightforward approach that keeps the $\beta$
fixed across all target samples and randomly draws from the errors'
empirical distribution of the training sample also seems to work well
in this scenario (OLS Empirical and Lasso Empirical). Perhaps the
worst performing approach is random forest (RF) exhibiting bias and
MSE that is considerably larger than that of other methods (Fig. \ref{fig:nonnormal_RF}).
This simulation also illustrates the importance of evaluating methods
across the entire welfare distribution, since if RF were evaluated
at the 50th percentile it would look be quite competitive with other
methods, when it considerably lags when evaluated at other thresholds.
These findings also carry over to the fully simulated data described
in section \ref{subsec:DGP-1--} and \ref{subsec:DGP-2--}, with the
most salient finding being that transformations of the data are of
little help in this instance and lead to larger a bias and MSE. Finally,
in these simulations the straightforward approach of keeping the $\beta$
fixed across all target samples and randomly drawing the residualsof
the training sample (OLS BS) outperforms other approaches (Figure
\ref{fig:sim_dgp_nonnormal_error}).

\subsubsection{Under heteroskedastic errors}

An additional simulation is conducted where the error distribution
of the DGP is modified to follow heteroskedastic errors using the
Mexican data as described in Section \ref{subsec:DGP-Mex}. The errors
are designed to mimic a situation where the uncertainty or variability
in the dependent variable grows as an independent variable increases.
Under these simulations, under the Mexican data, all methods work
relatively well with bias under 1 percentage point across all poverty
thresholds (Figure \ref{fig:bias_het}). However, the best performing
approach here is one that follows the alpha model from \citet{elbers2002micro}
leading to lower bias and MSE than the method \citet{harvey1976estimating}
but fit with MLE (``Alpha model'' versus ``MLE with het.'').\footnote{The method from \citet{harvey1976estimating} is applied via Stata's
\texttt{hetregress} command. While the alpha model from \citet{elbers2002micro}
is applied using \texttt{hetmireg (\citealp{corral_hetmireg_2018}).}} Nevertheless, in simulations following the DGP described in section
\ref{subsec:DGP-1--} and \ref{subsec:DGP-2--}, the MLE approach
by \citet{harvey1976estimating} seems to work better (Figure \ref{fig:sim_dgp_het}).
Ignoring the heteroskedasticity in the data, yet drawing residuals
from the empirical distribution (``OLS Empirical'' and ``Lasso
Empirical'') works better than just assuming the errors are homoskedastic
(``OLS ignoring het.''). This is aligned to the statement made by
\citet{betti2024attempt} who point out that ignoring heteroskedasticity
can lead to biased estimates. 

\begin{figure}[H]
\caption{Bias (left) and MSE (right) in FGT0 under different methods under
heteroskedasticity \label{fig:bias_het}}

\begin{centering}
\includegraphics[scale=0.62]{C:/Users/WB378870/GitHub/S2S_RTM/5.figures/3.Mexican/bias_het_mex}\includegraphics[scale=0.62]{C:/Users/WB378870/GitHub/S2S_RTM/5.figures/3.Mexican/mse_het_mex}
\par\end{centering}
{\footnotesize{}Note: Data are generated as described in \ref{subsec:Creating-populations}.
Bias is assessed at various poverty lines across the welfare distribution.
Specifically, these lines correspond to percentiles that are multiples
of 5. }\textbf{\footnotesize{}Fixed B}{\footnotesize{}: uses the method
described in Eq. \ref{eq:Fixed} to generate predictions. For each
of 100 Monte Carlo (MC) simulations, poverty is estimated at each
threshold and then averaged across simulations to produce the final
estimate. }\textbf{\footnotesize{}Fixed B lnskew}{\footnotesize{}:
follows the same approach as ``Fixed B'' but the dependent variable
is transformed using a zero-skewness log transformation. }\textbf{\footnotesize{}Fixed
B bcskew}{\footnotesize{}: follows the same approach as ``Fixed B''
but the dependent variable is transformed using a Box-Cox transformation.
}\textbf{\footnotesize{}MI X}{\footnotesize{}: implements X imputations
following Eq. \ref{eq:MI}, poverty is estimated for each threshold
under each imputed vector, and the results are averaged across simulations
to obtain the final estimate. }\textbf{\footnotesize{}MI X BS}{\footnotesize{}:
same as ``MI X'' but uses Stata's bootstrap option. }\textbf{\footnotesize{}lasso
BIC}{\footnotesize{}: applies parameters from lasso model where $\lambda$
is selected using the Bayesian information criterion (BIC) function
and residuals for each imputation are randomly drawn from the predicted
residuals. }\textbf{\footnotesize{}Het. MLE}{\footnotesize{}: uses
parameters from the \citet{harvey1976estimating} model for heteroskedasticity
and follows Eq. \ref{eq:Fixed} for the predictions where 100 MC simulations
are undertaken except that $\sigma_{e_{i}}^{2}$ is household specific.
}\textbf{\footnotesize{}Alpha model}{\footnotesize{}: obtains parameters
following \citet{elbers2002micro} model for heteroskedasticity and
follows Eq. \ref{eq:MI}, except that $\sigma_{e_{i}}^{2}$ is household
specific,with 100 imputed vectors.}{\footnotesize\par}
\end{figure}


\subsubsection{PMM - A flexible imputation approach }

A common alternative imputation approach, predictive mean matching
(PMM), differs in its simulation strategy. Rather than drawing residuals
or estimating parameters from a posterior distribution, PMM matches
observations from the source dataset to those in the target dataset
based on predicted values, ensuring that imputed values remain within
the observed range of the source data.\footnote{Predictions typically rely on OLS, see \citep{lucchetti2024application}
for an application to lasso model fitting.} This method is particularly useful when handling non-normal or discontinuous
distributions \citep{SwiftMethod}. However, its reliance on observed
values limits its ability to capture shifts in economic conditions
between periods, similar to other approaches presented in the section.
Nevertheless, recent applications that rely on imputing linking variables,
such as income-to-consumption ratios, are able to capture shifts in
economic conditions.\footnote{For a recent application see \citet{fernandez2024reconstructing}
and for a validation see \citet{stemmler2024beyond}. }

To illustrate PMM\textquoteright s limitations in its basic application,
a simulation where economic growth of 20 percent is introduced between
the source and target periods.\footnote{The only change introduced is a change in the constant term, everything
else remains equal between source and target data.} The simulation is designed to show how PMM behaves when the set of
values eligible for matching has less overlap with the true values.
Note that such an issue is likely to affect other methods as well,
and is not unique to PMM. In this scenario, PMM fails to capture the
distributional shift (Figure \ref{fig:bias_pmm-1}). This highlights
a critical limitation of PMM as well as other methods illustrated
in this paper: when substantial economic or structural changes occur
between data collection periods, the methods may not be appropriate
for imputing poverty estimates. This has also precipitated the innovations
introduced by \citet{fernandez2024reconstructing} where the PMM is
done on a linking variable and not directly to the dependent variable
of interest.

\begin{figure}[H]
\caption{Difference in FGT0 under PMM imputations \label{fig:bias_pmm-1}}

\begin{centering}
\includegraphics[scale=0.8]{C:/Users/WB378870/GitHub/S2S_RTM/5.figures/1.Normal/method_comp_PMMpp}
\par\end{centering}
{\footnotesize{}Note: Data are generated as described in \ref{subsec:Creating-populations}.
Difference is assessed at various poverty lines across the welfare
distribution of the source data. Specifically, these lines correspond
to percentiles that are multiples of 5. }\textbf{\footnotesize{}Predicted
change: }{\footnotesize{}applies PMM imputation using Stata's }\texttt{\footnotesize{}impute
pmm}{\footnotesize{} command.}{\footnotesize\par}
\end{figure}

These findings underscore that no single imputation method is universally
optimal. The choice of approach should be guided by the data\textquoteright s
properties, the error distribution, and the degree of economic and
social change between surveys. Parametric methods can introduce bias
when their assumptions are violated, while empirical residual drawing
and heteroskedasticity-aware models tend to offer more robust alternatives.
Machine learning techniques, particularly when combined with bootstrapping,
also show promise for future applications.

The simulation results demonstrate that while many imputation methods
perform well under ideal conditions, their effectiveness varies considerably
when faced with different error structures and data characteristics.
Direct application of parameters (\textquotedbl Fixed B\textquotedbl )
yields low bias under normal error distributions, but when assumptions
are violated---as with Student\textquoteright s t-distributed or
heteroskedastic errors---methods using empirical residual drawing
perform better. Machine learning approaches, such as lasso regression
combined with empirical residuals, show potential but require careful
implementation. Conversely, the limitations of Predictive Mean Matching
(PMM) become more pronounced when economic conditions change between
the training and target periods---though, such instability can affect
all models.

Overall, these results highlight the importance of selecting imputation
methods that are well-suited to the specific data context, error distribution,
and stability of conditions across surveys. Future work could usefully
extend bootstrap-based techniques to a wider range of machine learning
methods, while ensuring robust error estimation remains central.

\subsection{Imputations under biased samples}

Survey-to-survey (S2S) imputation methods are often applied to biased
samples, such as when specific population segments are systematically
underrepresented. To address this issue, practitioners commonly attempt
two corrective strategies: re-weighting the sample or standardizing
covariates. This section examines these methods and evaluates their
effectiveness using simulations.

\subsubsection{Extracting samples \label{subsec:Extracting-samples}}

Under each of the 500 populations created in Section \ref{subsec:DGP-1--}
and \ref{subsec:DGP-2--}, the following samples are taken: 
\begin{enumerate}
\item Random sample: a 20 percent simple random sample (SRS) of the population.
\item Bottom biased sample: here, the poorest quintile is purposely undersampled. 
\begin{enumerate}
\item For the top 80, take a SRS for each centile, $c$. 
\item For the bottom 20, take a $c$ percent as a sample. This means that
for the 20th centile an SRS sample of 20\% is taken, for the 19th
centile, 19 percent is sampled, for the 18th centile, 18 percent is
sampled, and so forth.
\end{enumerate}
\item Biased sample top and bottom: Create a biased sample, where the poorest
quintile and the richest quintile are purposely under sampled. 
\begin{enumerate}
\item For the top 20, sample $100-c$ percent. This means that for the 80th
centile, an SRS sample of 20 percent is taken, for the 81st centile
a sample of 19 percent is taken, for the 82nd centile, 18 percent
is sampled, and so forth. Note that we do not sample the top 1 percent. 
\item For the bottom 20, we sample $c$ percent. This means that for the
20th centile an SRS sample of 20 percent is taken, for the 19th centile,
19 percent is sampled, for the 18th centile, 18 percent is sampled,
and so forth. Note that the bottom 1 percent is not sampled. 
\item For all other centiles (21-79), an SRS sample by centile is taken. 
\end{enumerate}
\end{enumerate}

\subsubsection{Re-weighting the target sample}

Re-weighting involves adjusting the sampling weights in the target
survey to align its characteristics with those of the overall population.
One common approach, minimum cross-entropy,\footnote{When the original sampling weights of the target survey are ignored,
it is assumed that every observation had a similar probability of
being selected, then the re-weighting method is max-entropy (see Golan,
Judge and Miller \citeyearpar{GolanBook} for a more thorough exposition
of maximum entropy).} modifies the weights to match constraints like the mean or variance
of key covariates. While theoretically sound, re-weighting assumes
that adjusted weights can restore the representativeness of a biased
sample.

To test the problems with biased samples and re-weighting, two sampling
scenarios are used as an illustrative example. The samples are taken
from a population created as described in section \ref{subsec:Extracting-samples}.
For the biased sample exercise, three possible re-weighting scenarios
are considered:\footnote{Re-weighting is done using Stata's user created command \texttt{wentropy}
\citep{Corral_Rodas_wentropy}. The command is used instead of \citet{wittenberg2010introduction}
\texttt{maxentropy} command due to \texttt{wentropy} adding more flexibility
and because it provides solutions to instances where \texttt{maxentropy}
fails.} 
\begin{enumerate}
\item \textbf{Covariate match:} The mean of each covariate from the simple
random sample is used as a benchmark to adjust the weights of the
biased samples. Re-weighting is performed using minimum cross-entropy,
which modifies the prior weights as minimally as possible to ensure
that the adjusted sample meets the specified constraints, namely the
covariate means.
\item \textbf{Linear fit match:} The mean of the model's linear fit from
the simple random sample is used to adjust the weights of the biased
samples. This calibration is performed using minimum cross-entropy,
ensuring that the prior weights are modified as minimally as necessary
to satisfy the constraint of matching the linear fit's mean.
\item \textbf{Linear fit \& variance match:} The linear fit\textquoteright s
\textbf{mean} and \textbf{variance} of the simple random sample are
used to calibrate the weights of the biased samples. Weights are calibrated
using minimum cross-entropy to adjust prior weights by the smallest
amount possible so that the constraints (i.e., the mean and the variance
of the linear fit) are satisfied.\footnote{Since the simulated data does not have sampling weights, the priors
are equal to 1 and thus minimum cross-entropy is in reality max-entropy
in this instance.}
\end{enumerate}

\subsubsection{Standardizing covariates}

Standardizing covariates offers a straightforward approach to addressing
disparities between survey samples. Using data from the source survey
(or training data), where the model is initially estimated, practitioners
can leverage the known means and variances of each covariate to adjust
the corresponding variables in the target survey. This alignment ensures
that the mean and variance of covariates in the imputed data match
those of the source data, enhancing comparability. \citet{dang2017}
advocate this approach as a practical solution for reconciling differences
between survey samples. Their application of the method to Jordanian
data suggests that it can yield valid results under certain conditions.

The method's efficacy depends on the assumption that covariates follow
a normal distribution -- an assumption that may not hold universally.
To assess the robustness of this approach, the samples described in
Section \ref{subsec:Extracting-samples} were used, despite the fact
that the covariates in this case deviate from normality. The adjustment
process involves scaling and centering the covariates from the target
survey so that their means and variances align with those of the source
survey, as shown in the following formula:

\[
x_{target}^{'}=(x_{target}-\hat{\mu}_{x_{target}})\frac{\hat{\sigma}_{x}}{\hat{\sigma}_{x_{target}}}+\hat{\mu}_{x}
\]

Two standardization scenarios are implemented:
\begin{enumerate}
\item \textbf{Standardize each X}: The \textbf{mean} of each covariate from
the simple random sample is used as a benchmark to adjust the \textbf{mean
}and\textbf{ standard deviation} of the biased samples.
\item \textbf{Standardize XB}: The linear fit\textquoteright s ($\hat{y}=x\hat{\beta}$)
\textbf{mean} and \textbf{standard deviation} of the simple random
sample are used to standardize the linear fit of the biased target
samples. 
\end{enumerate}

\subsubsection{Results\label{subsec:Results_simdata}}

To isolate the impact of biased sampling in the target data and assess
the effectiveness of correction methods, the analysis avoids introducing
additional sources of error by using the true values for the coefficients
($\beta$) and error distribution ($\sigma_{e}^{2}$) applied in subection
\ref{subsec:DGP-1--} and \ref{subsec:DGP-2--}. Hence, instead of
fitting a model, the actual values of the linear fit from Eq. \ref{eq:dgp_model}
are used. For each sample created in Section \ref{subsec:Extracting-samples},
100 vectors were generated by combining the true $\beta$ with a randomly
drawn error term $e_{i}\sim N(0,0.5^{2})$ for every observation.\footnote{Note that the true value is used and not its estimate $\hat{\sigma}_{e}^{2}$.}
Poverty and Gini coefficients were then calculated for each of these
100 vectors using either re-weighted survey weights or standardized
covariates. The final estimates were obtained by averaging across
the 100 simulations, and bias was determined by comparing these estimates
to the true population values across the 1,000 populations, resulting
in a robust evaluation of the methods.

The most glaring outcome of the simulation is that adjustments to
correct for the sampling bias are unlikely to work when the linear
portion of the data (i.e., $x\beta$) are not normally distributed
(subsection \ref{subsec:DGP-2--}). Under this DGP none of the applied
adjustments led to improvements in bias (Figure \ref{fig:bias_reweight},
right). When the linear portion of the data are closer to normality
(subsection \ref{subsec:DGP-1--}), some of the adjustments may work
(Figure \ref{fig:bias_reweight}, left). The often applied approach
of reweighting the sample to match the covariate means to those of
the population (\textquotedblleft Match X\textquotedblright )\footnote{Since the weights assume no priors, these are in fact calculated using
maximum entropy. When prior weights are considered the method adjusts
the existing weights by the smallest possible amount to ensure the
constraints match, this is known as cross-entropy.} is of little use, even when the linear portion of the data is more
closely aligned to normality. Under this type of DGP, what does seem
to work is adjusting the entire linear fit ($x\beta$) to ensure its
mean and variance matches that of the target population (``Match
XB and var{[}XB{]}''). Similarly, standardizing the linear fit to
ensure its mean and variance are aligned to the target population's
also yields good results (``Standardize XB''), while standardizing
each covariate does not yield as good results. These methods effectively
address the interaction between mean and variance in determining poverty
estimates, but only seem to work when the DGP of the linear portion
of the model is close to normality. However, it is important to note
that the adjustments made here mimic the moments of the source data,
thus when applied to other data would just replicate the source data's
welfare distribution. Consequently the approaches would be of use
for contemporaneous S2S, but not for S2S over time.

\begin{figure}[H]
\caption{Bias in FGT0 under bottom biased samples (Sec \ref{subsec:Extracting-samples})
under different correction measures - DGP \ref{subsec:DGP-1--} (left)
and DGP \ref{subsec:DGP-2--} (right) \label{fig:bias_reweight}}

\begin{centering}
\includegraphics[scale=0.62]{C:/Users/WB378870/GitHub/S2S_RTM/5.figures/1.Normal/weights_std_bb_fgt0_bias}\includegraphics[scale=0.62]{C:/Users/WB378870/GitHub/S2S_RTM/5.figures/2.Nonnormal/weights_std_bb_fgt0_bias}
\par\end{centering}
{\footnotesize{}Note: Target samples are generated as described in
\ref{subsec:Extracting-samples}. Bias is assessed at various poverty
lines across the welfare distribution. Specifically, these lines correspond
to percentiles that are multiples of 5.}{\footnotesize\par}
\end{figure}

\begin{figure}[H]
\caption{Bias in FGT0 under top and bottom biased samples (Sec \ref{subsec:Extracting-samples})
under different correction measures - DGP \ref{subsec:DGP-1--} (left)
and DGP \ref{subsec:DGP-2--} (right) \label{fig:bias_reweight-1}}

\begin{centering}
\includegraphics[scale=0.62]{C:/Users/WB378870/GitHub/S2S_RTM/5.figures/1.Normal/weights_std_tbb_fgt0_mse}\includegraphics[scale=0.62]{C:/Users/WB378870/GitHub/S2S_RTM/5.figures/2.Nonnormal/weights_std_tbb_fgt0_mse}
\par\end{centering}
{\footnotesize{}Note: Target samples are generated as described in
\ref{subsec:Extracting-samples}. Bias is assessed at various poverty
lines across the welfare distribution. Specifically, these lines correspond
to percentiles that are multiples of 5.}{\footnotesize\par}
\end{figure}

For top-and-bottom-biased samples, where both the poorest and wealthiest
households were under-sampled, the performance of the correction methods
mirrored the earlier findings (Figure \ref{fig:bias_reweight-1}).
Unadjusted estimates again misrepresented poverty and inequality,
while methods that accounted for the mean and variance of the linear
fit provided more accurate results, but only for the DGP in subsection
\ref{subsec:DGP-1--}. Importantly, results under the DGP from subsection
\ref{subsec:DGP-2--} suggest that the methods should be applied with
caution and practitioners should recognize that these may not fully
address sample deviations from the population.

\textbf{Implications for Gini Estimation}

The challenges extend to inequality measurement. Biasing a sample
by excluding low-income households reduces $\mathrm{var}\left[x\beta\right]$,
which standard re-weighting methods fail to address (i.e. only addressing
differences in the mean). This limitation results in downward-biased
Gini estimates (Figure \ref{fig:bias_reweight-gini}). The issue is
even more pronounced for top-and-bottom-biased samples, where $\mathrm{var}\left[x\beta\right]$
is further diminished (Figure \ref{fig:bias_reweight-gini-1}). Adjusting
weights to align both the mean and variance of $x\beta$ with those
of the source data significantly improved the accuracy of the Gini
estimates. 

\begin{figure}[H]
\caption{Bias in Gini under bottom biased samples (Sec \ref{subsec:Extracting-samples})
and different correction measures - DGP \ref{subsec:DGP-1--} (left)
and DGP \ref{subsec:DGP-2--} (right) \label{fig:bias_reweight-gini}}

\begin{centering}
\includegraphics[scale=0.62]{C:/Users/WB378870/GitHub/S2S_RTM/5.figures/1.Normal/weights_std_bb_gini}\includegraphics[scale=0.62]{C:/Users/WB378870/GitHub/S2S_RTM/5.figures/2.Nonnormal/weights_std_bb_gini}
\par\end{centering}
{\footnotesize{}Note: Target samples are generated as described in
\ref{subsec:Extracting-samples}. Bias is assessed by comparing the
mean predicted Gini against the mean Gini of the populations generated
(i.e., the truth).}{\footnotesize\par}
\end{figure}

\begin{figure}[H]
\caption{Bias in Gini under top and bottom biased samples (Sec \ref{subsec:Extracting-samples})
and different correction measures - DGP \ref{subsec:DGP-1--} (left)
and DGP \ref{subsec:DGP-2--} (right) \label{fig:bias_reweight-gini-1}}

\begin{centering}
\includegraphics[scale=0.62]{C:/Users/WB378870/GitHub/S2S_RTM/5.figures/1.Normal/weights_std_tbb_gini}\includegraphics[scale=0.62]{C:/Users/WB378870/GitHub/S2S_RTM/5.figures/2.Nonnormal/weights_std_tbb_gini}
\par\end{centering}
{\footnotesize{}Note: Target samples are generated as described in
\ref{subsec:Extracting-samples}. Bias is assessed by comparing the
mean predicted Gini against the mean Gini of the populations generated
(i.e., the truth).}{\footnotesize\par}
\end{figure}


\subsection{Imputation over Time\label{subsec:over-time}}

Despite mildly encouraging results for correcting bias in cross-sectional
samples, these methods cannot be directly applied to S2S imputation
across time. The main limitation is the lack of reliable estimates
for the mean and variance of the dependent variable in the target
period. Without these benchmarks, standardization or reweighting becomes
ineffective, compounding the challenges of imputing poverty estimates
in dynamic contexts.

Moreover, when applying S2S to real-world data, the model used is
just one of many possible choices. The implicit assumption is that
this model represents the true data-generating process (DGP). In practice,
there are few clear criteria to select one model over another for
poverty imputation. Even a model that fits the source data well offers
little guarantee of robust performance when applied across time. Assuming
that the estimated model remains valid over time is a particularly
strong assumption. As noted by \citet{dang2025poverty}, a high model
$R^{2}$ does not necessarily translate into accurate poverty predictions.
As this section aims to show, multiple potential sources of bias can
arise when applying S2S over time.

Because household living standards surveys are costly and infrequently
collected---especially in low-income countries---S2S has become
an appealing option in many contexts. Its main applications over time
include::
\begin{enumerate}
\item Producing comparable poverty estimates: In many countries, surveys
may exist, but the welfare aggregate is not comparable across rounds.
S2S has been used to bridge this gap by training a model on the original
welfare aggregate and applying it to a more recent survey (e.g., Zambia
as noted by \citealt{yoshida2024enabling}). This process can also
be reversed---if the preferred welfare aggregate is from the newer
survey, the model can be trained on that data and applied retrospectively
(e.g., Nigeria, \citealt{lain2022estimating}).
\item Generating poverty estimates when no living standards data are available:
In many cases, a recent survey (e.g., a DHS or labor force survey)
does not include an adequate welfare aggregate. Here, a model trained
on the last available living standards survey is used to impute poverty
in the new survey (For example, see \citet{edochie2022we,lanjouw2024imputation,newhouse2019estimating}).
The target survey is often assumed to be nationally representative,
though this is not always true---sampling bias can distort results
(see \citet{sinha2022poverty}).
\end{enumerate}
S2S over time introduces additional assumptions beyond those outlined
in section \ref{subsec:sametime}. A key assumption when imputing
to a different period than the one used to train the model is that
the estimated parameters remain constant over time \citep{newhouse2014survey}.
This implies that the relationship between covariates and the dependent
variable is stable, and the distribution of unobservables does not
change. Essentially, this assumes that any shifts in the welfare distribution
are entirely driven by changes in the covariates and only the covariates
in the model--a strong assumption that becomes less likely to hold
the more time passes between source and target data.

Early work on S2S \citep{stifel2007tracking} often recommended selecting
covariates that may vary over time but maintain a stable relationship
with welfare -- a challenging criterion to meet in practice. Furthermore,
restricting the eligible covariates for the model can result in insufficient
explanation of the variation in welfare. This, in turn, increases
reliance on the distribution of unobservables, which is based on the
source data used for the model and may not reflect the target period
accurately. 

In this section the discussion focuses on how changes in the parameters
can affect the poverty predictions. Through simulated data, components
can be adjusted one at a time to identify how these may impact predictions.
In practice, it is possible that any of the parameters used for the
imputation $\left(\beta,\:\sigma\right)$ has changed over time, however
it is impossible to know if the changes compound biases or cancel
each other. This requires transparent communication and careful consideration
when applying these methods and so potential users, including policy
makers, fully comprehend the limitations of the predictions.

\subsubsection{Changes in the Error's Distribution}

It is possible, though unlikely, for the relationship between the
covariates and the dependent variable to remain constant over time
while poverty levels still change. The change can be driven entirely
by changes in the distribution of the unobservables, $\sigma_{e}^{2}$.
Changes in the unobservables also implies a change in inequality.
To illustrate this, the same population created in sections \ref{subsec:DGP-1--}
and \ref{subsec:DGP-2--} are used. The simulation consists in keeping
the linear fit constant and only change the value of $\sigma_{e}$.
This is done across the 1,000 populations where predictions are obtained
from Eq. \ref{eq:Fixed}.

Results for the simulation are presented in Figure \ref{fig:Change-in-poverty_sigma},
and illustrate how poverty rates change under different lines determined
at the percentile of the original welfare. Hence, if the original
poverty rate was 20 percent, and using that same corresponding poverty
line when increasing the value of $\sigma_{e}$ by 20 percent, the
resulting poverty rate would be biased upwards by over 2 percentage
points (Figure \ref{fig:Change-in-poverty_sigma}, left). The bias
could be greater or smaller and depends on the DGP, under the DGP
from section \ref{subsec:DGP-2--}, the bias is even higher. At higher
thresholds the gap between the original poverty rate and the new rates
is less noticeable, although still present. Notably, estimates at
the 50th percentile show no bias, reinforcing the importance of evaluating
poverty predictions across the full welfare distribution.

\begin{figure}[H]
\caption{Change in poverty prediction if $\sigma$ changes by $x\%$ - DGP
\ref{subsec:DGP-1--} (left) and DGP \ref{subsec:DGP-2--} (right)\label{fig:Change-in-poverty_sigma}}

\begin{centering}
\includegraphics[scale=0.62]{C:/Users/WB378870/GitHub/S2S_RTM/5.figures/1.Normal/sigma_change_fgt0}\includegraphics[scale=0.62]{C:/Users/WB378870/GitHub/S2S_RTM/5.figures/2.Nonnormal/sigma_change_fgt0}
\par\end{centering}
{\footnotesize{}Note: Target samples are generated as described in
\ref{subsec:Extracting-samples}, only the SRS samples are used for
this simulation. Bias is assessed at various poverty lines across
the welfare distribution. Specifically, these lines correspond to
percentiles that are multiples of 5. Predictions are obtained following
Eq. \ref{eq:Fixed}, where the value for $\sigma_{e}^{2}$ is changed.}{\footnotesize\par}
\end{figure}


\subsubsection{Changes in the Constant Term}

When conducting imputations over time, changes to the constant term
are seldom considered. Failing to capture a change in the constant
term can lead to relatively stagnant predictions over time, leading
to predictions that are considerably off the mark. In Figure \ref{fig:Change-in-poverty_intercept}
the true constant term under DGP \ref{subsec:DGP-1--} (left) and
DGP \ref{subsec:DGP-2--} (right) is adjusted by $x\%$ in the data
generating process for the target data. Even slight changes lead to
considerable differences in poverty. Imagine a scenario where transformed
welfare for everyone has increased by $1\%$ (Under GDP growth, for
example.), yet everything else about the welfare distribution remains
the same. This would entail a neutral distribution shift to the right,
and thus everyone's transformed welfare is improved by 1\%. However,
relying on a model fit on the source data, before the increase --
even if everything else in the model remains the same -- we would
essentially ignore the growth in welfare. This would lead to imputations
that considerably overestimate poverty for the new period.

\begin{figure}[H]
\caption{Resulting poverty prediction if constant term changes by $x\%$ -
DGP \ref{subsec:DGP-1--} (left) and DGP \ref{subsec:DGP-2--} (right)
\label{fig:Change-in-poverty_intercept}}

\begin{centering}
\includegraphics[scale=0.62]{C:/Users/WB378870/GitHub/S2S_RTM/5.figures/1.Normal/constant_change_fgt0}\includegraphics[scale=0.62]{C:/Users/WB378870/GitHub/S2S_RTM/5.figures/2.Nonnormal/constant_change_fgt0}
\par\end{centering}
{\small{}Source: Based on simulated data illustrated in Section \ref{subsec:Creating-populations}.
Bias is assessed at various poverty lines across the welfare distribution.
Specifically, these lines correspond to percentiles that are multiples
of 5. }{\footnotesize{}Predictions are obtained following Eq. \ref{eq:Fixed}.}{\footnotesize\par}
\end{figure}


\subsubsection{Omitted Variables \label{subsec:Omitted-Variables}}

Omitted variables and endogeneity have traditionally not been major
concerns in predictions such as S2S or small area estimation. In a
contemporaneous setting, OLS fitted values remain unbiased for the
conditional mean under standard conditions, with unobserved determinants
absorbed---on average---by the intercept and by coefficients on
correlated regressors. This reassurance is weaker in applications
over time: if the distribution of omitted determinants changes between
the source and target periods, the source-period relationship between
observables and welfare need not remain transferable, and prediction
errors can become systematic. For a more detailed discussion, see
the annex (Annex \ref{subsec:ann_ovb}). 

Consider a plausible application of survey-to-survey (S2S) imputation
for monetary poverty in a conflict-affected rural setting. Suppose
conflict declines between two survey rounds and, as a result, poverty
falls. If the imputation model does not control for conflict directly,
it may rely on fast-changing consumption indicators---such as whether
households consumed meat, eggs, or chocolate---as proxies for welfare
dynamics, consistent with guidance in \citet{yoshida2022concept}.
However, if conflict simultaneously depresses consumption of these
goods and reduces household expenditure, then omitting conflict from
the model induces omitted-variable bias: the coefficients on the consumption
indicators partly reflect the conflict channel, and changes in conflict
over time can translate into biased predictions of welfare and poverty
(refer to Annex \ref{subsec:ann_ovb}).

Proxy variables can, in principle, mitigate omitted-variable bias
\citep{wooldridge2009introductory,wooldridge2010econometric,YoshidaKawashimaTakamatsu2025ShockResponsiveness},
but not every correlated covariate qualifies as a valid proxy. \citet{wooldridge2009introductory},
for example, discusses \textquotedblleft ability\textquotedblright{}
as an omitted determinant of wages and notes that while education
is correlated with ability, it is itself a structurally relevant regressor
in the wage equation and therefore should not be treated as a proxy;
he instead uses IQ as a proxy for ability. Formally, a valid proxy
must satisfy two conditions \citep{wooldridge2009introductory,wooldridge2010econometric}:
(i) it must be strongly correlated with the omitted factor, and (ii)
it must be conditionally irrelevant for the outcome once the omitted
factor is included---i.e., if the omitted variable were observed,
the proxy would add no explanatory power to the model. Item-specific
consumption indicators may plausibly satisfy the first condition,
but they are unlikely to satisfy the second because they are mechanically
related to household expenditure, and thus remain directly relevant
determinants of the welfare aggregate. In other words, a valid proxy
should be related to expenditures only through its relationship with
the omitted variable. When this \textquotedblleft conditional irrelevance\textquotedblright{}
requirement fails, the proposed proxy is simply another regressor
that can inherit bias when it is correlated with the omitted factor,
potentially amplifying---rather than alleviating---prediction error.

To simulate the potential impact of omitted variable bias on poverty
predictions over time the data simulated in sections \ref{subsec:DGP-1--}
and \ref{subsec:DGP-2--} are expanded to include 2 new covariates: 
\begin{enumerate}
\item Conflict. Conflict is assumed to be negatively correlated with welfare
with a coefficient equal to $-0.3$. 
\begin{enumerate}
\item It is simulated as a binary variable taking a value of 1 when a random
uniform number between 0 and 1 is less than $c=0.4$. 
\end{enumerate}
\item Eggs purchased. The variable is assumed to be positively correlated
to welfare with a coefficient equal to $0.4$. However, the share
of households purchasing eggs can change due to conflict and external
factors unrelated to welfare.
\begin{enumerate}
\item The variable is simulated as binary, taking the value 1 if a randomly
drawn uniform number between 0 and 1 falls below $0.5x-0.5\times Conflict$,
where $x=1$ in the baseline. Importantly, \textquotedblleft eggs
purchased\textquotedblright{} is not generated from a behavioral demand
model. It is introduced as a reduced-form indicator whose probability
shifts with conflict and with an exogenous component $x$ that captures
non-welfare related drivers of purchase decisions---such as market
access and availability, supply disruptions, relative price changes,
or shifts in preferences. The objective is to mimic a setting in which
an item-level consumption indicator moves because the economic environment
changes along with changes in conflict.
\end{enumerate}
\end{enumerate}
Under the baseline simulation, the value for $x=1$ and $c=0.4$,
thus, roughly 40 percent of households are affected by conflict and
roughly 30 percent of households purchased eggs. Welfare is simulated
assuming the true DGP, and thus includes conflict and eggs with their
true coefficients. The prediction model is produced omitting conflict
and is executed using the complete data to avoid other potential sources
of bias (Table \ref{tab:Omitted-variable-model}). As presented in
Annex \ref{subsec:ann_ovb}, the coefficient on eggs is upward biased.
Due to the omission of conflict, the constant term of the regression
is downward biased. The adjustment in the constant term ensures that
the model's prediction of the dependent variable in the model data
is unbiased. As can be seen, omitting conflict from the model also
leads to a larger RMSE than the truth since it absorbs the explained
portion of the omitted variable, conflict, which will likely affect
poverty predictions over time. Nevertheless, for predictions to the
same period OVB has little if any impact (Figure \ref{fig:ovb same period}).\footnote{The reason for this is that even in the presence of OVB $\mathrm{var}\left[x\hat{\beta}\right]+\hat{\sigma}_{e}^{2}$
will be the same as the one where conflict is present and since $\hat{y}$
is also equal poverty predictions in the same period are unaffected.}

\begin{table}[H]
\caption{Omitted variable model comparison\label{tab:Omitted-variable-model}}

\centering{}%
\begin{tabular}{ccccc}
 & \multicolumn{2}{c}{DGP \ref{subsec:DGP-1--}} & \multicolumn{2}{c}{DGP \ref{subsec:DGP-2--}}\tabularnewline
\hline 
 & OVB Model & Complete Model & OVB Model & Complete Model\tabularnewline
\hline 
x1 & 0.0989 & 0.0989 & 0.0995 & 0.0993\tabularnewline
x2 & 0.488 & 0.490 & 0.498 & 0.494\tabularnewline
x3 & -0.247 & -0.244 & -0.247 & -0.246\tabularnewline
x4 & -0.201 & -0.201 & -0.211 & -0.210\tabularnewline
x5 & -0.160 & -0.157 & -0.155 & -0.149\tabularnewline
eggs & 0.577 & 0.402 & 0.584 & 0.411\tabularnewline
conflict &  & -0.307 &  & -0.299\tabularnewline
Intercept & 2.838 & 3.012 & 2.828 & 3.002\tabularnewline
\hline 
R2 & 0.536 & 0.565 & 0.581 & 0.644\tabularnewline
RMSE & 0.509 & 0.492 & 0.320 & 0.295\tabularnewline
Observations & 20,000 & 20,000 & 20,000 & 20,000\tabularnewline
$\hat{y}$ & 2.976 & 2.976 &  & \tabularnewline
\hline 
\end{tabular}
\end{table}

\begin{figure}[H]
\caption{Resulting poverty prediction under OVB in the same period \label{fig:ovb same period}}

\begin{centering}
\includegraphics[scale=0.4]{C:/Users/WB378870/GitHub/S2S_RTM/5.figures/1.normal/obv_pred_orig}\includegraphics[scale=0.4]{C:/Users/WB378870/GitHub/S2S_RTM/5.figures/2.nonnormal/obv_pred_orig}
\par\end{centering}
{\small{}Source: Based on simulated data illustrated in Section \ref{subsec:Creating-populations}
with added covariates. Predictions are assessed at various poverty
lines across the welfare distribution. Specifically, these lines correspond
to percentiles that are multiples of 5.}{\small\par}
\end{figure}

Once conflict subsides, so that the share of households exposed to
conflict is only 20 percent, the share of households purchasing eggs
will increase to 40 percent due to the assumed relationship. Note
that $x=1$in this scenario, thus there is no external shock to the
egg purchase decision beyond conflict. Using the adjusted covariates,
the true resulting welfare is calculated -- including conflict and
the correct coefficients. Then the model parameters obtained from
the original data, excluding conflict, are used to obtain predictions
of poverty across the welfare distribution. As can be seen, even in
the absence of any other changes to the DGP, the predictions under
the model will be considerably biased only due to a change in conflict
(Figure \ref{fig:ovb_Dconflict_noDegg}). 

\begin{figure}[H]
\caption{Resulting poverty prediction under OVB in a different period, under
DGP \ref{subsec:DGP-1--} (left) and DGP \ref{subsec:DGP-2--} (right)
\label{fig:ovb_Dconflict_noDegg}}

\begin{centering}
\includegraphics[scale=0.62]{C:/Users/WB378870/GitHub/S2S_RTM/5.figures/1.normal/bd5_ed10}\includegraphics[scale=0.62]{C:/Users/WB378870/GitHub/S2S_RTM/5.figures/2.nonnormal/bd5_ed10}
\par\end{centering}
{\small{}Source: Based on simulated data illustrated in Section \ref{subsec:Creating-populations}
with added covariates. Predictions are assessed at various poverty
lines across the welfare distribution. Specifically, these lines correspond
to percentiles that are multiples of 5. In this prediction, the share
of households exposed to conflict has been halved.}{\small\par}
\end{figure}

If there is any type of exogenous shock that would affect the household's
decision to purchase eggs, beyond just a change in conflict, then
biases could be much more pronounced or not meaningful. The direction
and magnitude of the bias is dependent on the change and differ depending
on the DGP used for the model. In some instances the omission is not
too problematic (Figure \ref{fig:ovb new period-2}), but in some
it is possible that predictions and what actually happened may be
completely misaligned (Figure \ref{fig:ovb new period-2-1}). This
is in agreement with the discussion in the Annex \ref{subsec:ann_ovb},
the direction of the prediction bias is not immediately clear. In
many instances the predictions will suggest an increase in poverty
whereas in reality poverty has dropped. In other instances it may
underestimate the total change. 

\begin{figure}[H]
\caption{Resulting poverty prediction under OVB in a different period \label{fig:ovb new period-2}}

\begin{centering}
\includegraphics[scale=0.62]{C:/Users/WB378870/GitHub/S2S_RTM/5.figures/1.normal/bd5_ed13}\includegraphics[scale=0.62]{C:/Users/WB378870/GitHub/S2S_RTM/5.figures/2.nonnormal/bd5_ed13}
\par\end{centering}
{\small{}Source: Based on simulated data illustrated in Section \ref{subsec:Creating-populations}
with added covariates. Predictions are assessed at various poverty
lines across the welfare distribution. Specifically, these lines correspond
to percentiles that are multiples of 5. In these simulations conflict
has halved and $x=1.3$, thus the share of households who purchase
eggs increased to 55 percent.}{\small\par}
\end{figure}

\begin{figure}[H]
\caption{Resulting poverty prediction under OVB in a different period \label{fig:ovb new period-2-1}}

\begin{centering}
\includegraphics[scale=0.62]{C:/Users/WB378870/GitHub/S2S_RTM/5.figures/1.normal/bd5_ed7}\includegraphics[scale=0.62]{C:/Users/WB378870/GitHub/S2S_RTM/5.figures/2.nonnormal/bd5_ed7}
\par\end{centering}
{\small{}Source: Based on simulated data illustrated in Section \ref{subsec:Creating-populations}
with added covariates. Predictions are assessed at various poverty
lines across the welfare distribution. Specifically, these lines correspond
to percentiles that are multiples of 5. In these simulations conflict
has halved and $x=1.3$, thus the share of households who purchase
eggs is 25 percent.}{\small\par}
\end{figure}

The impact of omitted variable bias extends beyond \textquotedblleft fast-changing\textquotedblright{}
consumption variables. Consider a model that includes an indicator
for subsistence farmers, who are more likely to be poor and thus negatively
correlated with welfare. Now imagine a massive drought occurs between
the year the model was trained and the target year. This drought could
reduce the number of people farming. If those former farmers are unable
to meet their needs, the model might underestimate poverty because
it omitted the effect of the drought. Another example involves the
omission of remittances which are positively related to welfare and
likely positively related to ``fast-changing'' consumption variables,
such as eggs or meat. Under this instance the coefficient on eggs
would be upward biased, and since the impact of the omitted variable
will be absorbed by the intercept it will likely also be upward biased.
Thus, changes in remittances that are not captured by the model are
likely to lead to biased predictions. 

The magnitude and direction of the bias that arises due to OVB is
often not clear and could compound other biases just as easily as
it could be offset by other biases. Nevertheless, unless actual welfare
data are available for the imputed year, it is impossible to truly
know how the model is performing. 

\section{Conclusions}

This paper has highlighted both the potential and the challenges of
survey-to-survey (S2S) imputation for poverty prediction, emphasizing
the need for cautious application and robust methodological development.
While S2S has several limitations, it may be the best available option
in certain scenarios -- particularly when real-time poverty estimates
are needed, but household surveys are outdated, unavailable, or inconsistent,
and when methods relying on real GDP growth are not feasible due to
data limitations or reliability concerns. In some cases, S2S may offer
more reliable estimates than alternatives such as extrapolations based
purely on national accounts, which often diverge from household survey-based
welfare measures. However, the likelihood of getting the numbers wrong
is not negligible and the main problem is that it is difficult to
assertain when the numbers are likely wrong.

Based on evidence from the simulations undertaken here, the following
recommendations emerge: 
\begin{enumerate}
\item \textbf{Understand and Address Bias Sources:} Practitioners should
assess changes in error distributions and parameter instability, as
these are key sources of prediction bias. This is especially critical
during periods of economic transition. Fast-moving predictors should
not be relied upon without thorough validation to mitigate omitted
variable bias and other factors affecting the distribution of unobservables. 
\item \textbf{Enhance Validation Practices:} Validation efforts should go
beyond same-period tests and include cross-time validation to assess
robustness against economic and structural changes. Practitioners
should also evaluate predictions across multiple poverty thresholds
to capture distributional nuances. Additionally, external validation
-- comparing results with broader indicators such as GDP and unemployment,
while incorporating country-specific knowledge -- enhances reliability.
Good examples of these can be seen in \citet{sinha2022poverty} as
well as \citet{lain2022estimating} who rely on auxiliary data to
lend further credence to their S2S estimates.
\item \textbf{Incorporate Structural and Temporal Dynamics:} Methods that
allow for time-varying parameters and account for structural changes
in predictor-welfare relationships should be considered. For instance,
shifting the constant term by observed growth between survey periods
could help adjust for evolving dynamics. Combining multiple data sources
may further improve the robustness and flexibility of predictions. 
\item \textbf{Foster Transparency in Methodology and Communication: }Practitioners
must pursue clear and transparent reporting of assumptions, validation
efforts, and uncertainties in S2S applications. Practitioners should
thoroughly document how imputations are implemented---how errors
and parameters are applied---and clearly communicate the limitations
of the method, especially in contexts with high economic volatility.
Uncertainty intervals should go beyond simple noise estimates and
incorporate scenario analyses to contextualize findings for stakeholders.
This may include sensitivity analyses that adjust key parameters to
illustrate how results might shift under different assumptions.
\item \textbf{Exercise Caution in Interpretation:} S2S methods should be
applied with explicit acknowledgment of their limitations, particularly
for tracking poverty over long periods. Transparency is essential
-- clearly stating that the estimates are predicted rather than derived
from an actual household survey ensures proper interpretation and
avoids misrepresentation. 
\item \textbf{Invest in Data and Research:} Greater investment in regular,
high-quality data collection is essential for reducing over-reliance
on imputation methods and for improving the accuracy of poverty measurement.
S2S methods ultimately depend on the quality and representativeness
of the data they are built upon---weak or outdated data will naturally
compromise the quality of imputed estimates. As the costs of data
collection rise, there is a growing temptation to substitute imputation
for new data. However, imputation should complement---not replace---regular
survey programs. Sustained investment in data is necessary to ensure
that models remain valid and are regularly recalibrated to reflect
real-world changes. This point is also highlighted in \citet{fujii2020predicted},
who caution against replacing genuine data with modeled estimates
in contexts where underlying conditions are evolving. 
\item \textbf{Support Survey Improvements While Maintaining Comparability}:
Statistical Offices should be encouraged to update their surveys in
line with best practices -- for example, refining the measurement
of housing and durable goods in consumption aggregates. However, such
improvements can compromise comparability with earlier surveys. To
address this, Statistical Offices could introduce \textquotedbl bridging\textquotedbl{}
surveys by fielding the old questionnaire to a sub-sample of the new
survey sample. This approach would allow them to assess key S2S assumptions,
such as changes in the constant term of prediction models, helping
to improve the reliability of imputed estimates.
\end{enumerate}
By recognizing both the challenges and the practical value of S2S,
we can better position it as a useful tool for addressing poverty
data gaps while maintaining the rigor needed for informed policy decisions.
S2S should not be seen as an automatic substitute for household surveys
but as a tool that -- when carefully applied and validated -- can
provide meaningful insights when other data sources are unavailable
or unreliable. The appropriate use of S2S is not just about improving
estimates but also about demonstrating when cruder alternatives may
be far more misleading. 

\pagebreak{}

\printbibliography

\pagebreak{}

\section{Annex}

\subsection{The Case of Omitted Variable Bias\label{subsec:ann_ovb}}

Survey-to-survey imputation across time periods requires careful consideration.
For example in a conflict setting, where conflict has significantly
subsided between the collection of the training and target data, it
is crucial to control for conflict in the model.\footnote{Note that nothing to control for those shocks was included in the
original model.} In instances where conflict is not included in the model, its effects
would only manifest through changes in other variables. This potential
omitted variable bias raises concerns about the model's reliability
for temporal predictions. When imputing over time, it is possible
for multiple issues to arise. Since true expenditure is not available
it is impossible to know if the model's biases cancel each other out
or compound.

Some may argue that the inclusion of fast changing consumption variables,
such as an indicator of whether or not the household consumed eggs
over the past week, is sufficient for accurate estimates of poverty
. Nevertheless, these variables are likely subject to omitted variable
bias which can lead to biased estimates of the coefficients leading
to biased poverty measures. In the case of the imputation example
given here, the omitted variable is conflict. In others it may be
the introduction of a cash transfer or a global pandemic, for example. 

Assuming conflict is negatively related to expenditure and negatively
related to consumption of certain goods, such as meat and eggs, this
would lead to coefficients that are upward biased. Assume the following,
simplified, model:

\begin{equation}
Y=\beta_{0}+\beta_{1}Eggs+\beta_{2}Conflict+\varepsilon\label{eq:full_egg_conflict}
\end{equation}

By omitting conflict the model is now:

\begin{equation}
Y=\tilde{\beta}_{0}+\tilde{\beta}_{1}Eggs+u\label{eq:fit_egg_conflict}
\end{equation}

where $u=\beta_{2}Conflict+\varepsilon$. In a model where conflict
is not included $\tilde{\beta}_{1}$would be greater than $\beta_{1}$.
The OLS estimator for $\tilde{\beta}_{1}$ is given by:

\[
\tilde{\beta}_{1}=\frac{\mathrm{Cov}\left(Eggs,Y\right)}{\mathrm{Var}\left[Eggs\right]}
\]

\[
\tilde{\beta}_{1}=\frac{\mathrm{Cov}\left(Eggs,\beta_{0}+\beta_{1}Eggs+\beta_{2}Conflict+\varepsilon\right)}{\mathrm{Var}\left[Eggs\right]}
\]

and relying on the linearity of covariance:

\[
\tilde{\beta}_{1}=\beta_{1}+\frac{\beta_{2}\mathrm{Cov}\left(Eggs,Conflict\right)}{\mathrm{Var}\left[Eggs\right]}
\]

Since we assume that conflict and welfare are negatively correlated,
and negatively correlated with the likelihood of buying eggs, then
we know that $\beta_{2}<0$ and that $\mathrm{Cov}\left(Eggs,Conflict\right)<0$.
This would lead to $\frac{\beta_{2}\mathrm{Cov}\left(Eggs,Conflict\right)}{\mathrm{Var}\left[Eggs\right]}>0$
which means that $\tilde{\beta}_{1}$ is upward biased. 

In addition, the intercept will be underestimated. The intercept is
equal to:

\[
\tilde{\beta}_{0}=\bar{y}-\tilde{\beta}_{1}\overline{Eggs}
\]

\[
\tilde{\beta}_{0}=\beta_{0}+\beta_{1}\overline{Eggs}+\beta_{2}\overline{Conflict}-\tilde{\beta}_{1}\overline{Eggs}
\]

\begin{equation}
\tilde{\beta}_{0}=\beta_{0}+\beta_{2}\overline{Conflict}+(\beta_{1}-\tilde{\beta}_{1})\overline{Eggs}\label{eq:intercept}
\end{equation}

we know that $\beta_{2}<0,$ and that $(\beta_{1}-\tilde{\beta}_{1})<0.$
This suggests that the intercept term is downward biased. 

When predicting over time, the direction of the prediction bias is
dictated by changes in conflict and the share of households who purchased
eggs. Any change in conflict, will lead to a change in egg purchases
and expenditure. The direction of the bias in prediction is undetermined:

\[
E\left[\bar{\tilde{y}}-\bar{y}\right]=E\left[\tilde{\beta}_{0}+\tilde{\beta}_{1}\overline{Eggs}'-\beta_{0}-\beta_{1}\overline{Eggs}'-\beta_{2}\overline{Conflict'}\right]
\]

replace $\tilde{\beta}_{0}$ with the value of Eq. \ref{eq:intercept},
and re-arrange:

\[
E\left[\bar{\tilde{y}}-\bar{y}\right]=E\left[\beta_{2}\left(\overline{Conflict}-\overline{Conflict'}\right)+\left(\beta_{1}-\tilde{\beta}_{1}\right)\left(\overline{Eggs}-\overline{Eggs}'\right)\right]
\]

then a prediction will be \textbf{upward biased} if:

\[
\beta_{2}\left(\overline{Conflict}-\overline{Conflict'}\right)>\left(\tilde{\beta}_{1}-\beta_{1}\right)\left(\overline{Eggs}-\overline{Eggs}'\right)
\]

and downward biased if:

\[
\beta_{2}\left(\overline{Conflict}-\overline{Conflict'}\right)<\left(\tilde{\beta}_{1}-\beta_{1}\right)\left(\overline{Eggs}-\overline{Eggs}'\right)
\]

and while they could cancel eachother out, it is unlikely.

Note that since poverty also depends on the model's predicted root
mean squared error ($RMSE$, $\hat{\sigma}$), which would change
over time since it is a function of the omitted variable -- conflict
-- then there are 2 sources of bias. The first is due to the egg
coefficient and the intercept, and the second is due to the difference
in the distribution of errors which changes due to the omitted variable.
A decrease in conflict would potentially lead to a shrinking of the
$RMSE$, lowering the poverty likelihood of households but would be
ignored in the example provided here.


\subsection{Data generation under a two-fold nested error model\label{subsec:DGP_cluster}}

Data is generated for the simulations following the assumed double
nested error model.\footnote{Data is simulated following the same approach as \citet{corral2022guidelines}
which adapts the DGP from \citet{ea7f7e4f-7381-39b0-abdd-8da14c150408}.
The write-up here is also borrowed from \citet{corral2022guidelines}.} The simulated population consists of $N=20,000$, distributed across
$A=40$ areas $\left(a=1,\ldots,A\right)$. Within each area $a$,
observations are uniformly allocated over $c=10$ clusters $\left(c_{a}=1,\ldots,C_{a}\right)$.
Each cluster $c$ consists of $N_{ac}=50$ observations and each area
consists of $N_{a}=500$ observations.

The outcome variable $Y$ is generated from a linear model with additive
random effects at the area and PSU level, and household-specific error:

\begin{equation}
\ln y_{aci}=3+0.1x_{1_{i}}+0.5x_{2_{i}}-0.25x_{3_{i}}+0.2x_{4_{i}}-0.15x_{5_{i}}+\eta_{a}+\eta_{ac}+\varepsilon_{aci}\label{eq:dgp_model-1}
\end{equation}

\begin{enumerate}
\item $x_{1}$ is a binary variable, taking value 1 when a random uniform
number between 0 and 1, at the household-level, is less than or equal
to $0.3+0.5\frac{a}{40}+0.2\frac{c}{10}$. 
\item $x_{2}$ is a binary variable, taking value 1 when a random uniform
number between 0 and 1 is less than 0.2.
\item $x_{3}$ is a binary variable, taking value 1 when a random uniform
number between 0 and 1 is less than 0.5 as long as $x_{2}=1$, otherwise
it is equal to 0.
\item $x_{4}$ is a binary variable, taking value 1 when a random uniform
number between 0 and 1, at the household-level, is less than or equal
to $0.5+0.3\frac{a}{40}+0.1\frac{c}{10}$. 
\item $x_{5}$ is a variable drawn from a Student\textquoteright s t distribution
with 5 degrees of freedom and scaled by 0.25.
\end{enumerate}
Random effects are generated as follows:
\begin{itemize}
\item Cluster effects are simulated as $\eta_{ac}\stackrel{iid}{\sim}N\left(0,0.1\right)$
\item Area effects as $\eta_{a}\stackrel{iid}{\sim}N\left(0,0.15^{2}\right)$
\item Household specific residuals as $e_{aci}\overset{iid}{{\sim}}N\left(0,0.5^{2}\right)$,
where $c=1,\ldots,C_{a}$, $a=1,\ldots,A$. 
\end{itemize}
For each of the 1,000 simulated populations, a two-stage sampling
process is used to draw the source sample (used to train the model):
\begin{enumerate}
\item In the first stage, 4 clusters are selected at random (without replacement)
within each area.
\item In the second stage, 10 observations are randomly selected from each
selected cluster.
\end{enumerate}
This results in a sample of 1,600 observations per population. The
target data is drawn using the same sampling procedure, so the target
sample may include entirely different clusters and observations than
the source sample.

\subsection{Additional figures}

\begin{figure}[H]
\caption{Bias and MSE in FGT0 under a two-fold double nested normal error DGP
\label{fig:bias_mse_cluster}}

\begin{centering}
\subfloat[Bias under DGP \ref{subsec:DGP-1--} (left) and DGP \ref{subsec:DGP-2--}]{\begin{centering}
\includegraphics[scale=0.62]{C:/Users/WB378870/GitHub/S2S_RTM/5.figures/1.Normal/mi_ebp_fgt0_cluster_bias}\includegraphics[scale=0.62]{C:/Users/WB378870/GitHub/S2S_RTM/5.figures/2.Nonnormal/mi_ebp_fgt0_cluster_bias}
\par\end{centering}
}
\par\end{centering}
\begin{centering}
\subfloat[MSE under DGP \ref{subsec:DGP-1--} (left) and DGP \ref{subsec:DGP-2--}]{\begin{centering}
\includegraphics[scale=0.62]{C:/Users/WB378870/GitHub/S2S_RTM/5.figures/1.Normal/mi_ebp_fgt0_cluster_mse}\includegraphics[scale=0.62]{C:/Users/WB378870/GitHub/S2S_RTM/5.figures/2.Nonnormal/mi_ebp_fgt0_cluster_mse}
\par\end{centering}
}
\par\end{centering}
{\footnotesize{}Note: Data are generated using Mexican data as described
in \ref{subsec:Creating-populations}. Bias is assessed at various
poverty lines across the welfare distribution. Specifically, these
lines correspond to percentiles that are multiples of 5. }\textbf{\footnotesize{}OLS}{\footnotesize{}:
uses the method described in Eq. \ref{eq:Fixed} to generate predictions,
ignoring the nested error structure. }\textbf{\footnotesize{}Nested:
}{\footnotesize{}similar to OLS but considers double nested error
structure of the DGP. }\textbf{\footnotesize{}MI X}{\footnotesize{}:
implements X imputations following Eq. \ref{eq:MI}. The model is
trained on one sample taken from the Mexican data, and applied to
500 different samples taken from the data. Poverty thresholds are
determined using the census.}{\footnotesize\par}
\end{figure}

\begin{figure}[H]
\caption{Bias (left) and MSE (right) in FGT0 of different methods under non-normal
errors for the Mexican data \label{fig:nonnormal_RF}}

\begin{centering}
\includegraphics[scale=0.62]{C:/Users/WB378870/GitHub/S2S_RTM/5.figures/3.Mexican/bias_nonnormal8_mex_RF}\includegraphics[scale=0.62]{C:/Users/WB378870/GitHub/S2S_RTM/5.figures/3.Mexican/mse_nonnormal8_mex_RF}
\par\end{centering}
{\footnotesize{}Note: Data are generated as described in \ref{subsec:DGP-Mex};
errors are simulated from a Student's t-distribution with 8 degrees
of freedom and scaled for a SD of 0.8. Bias and MSE are assessed as
noted in Section\ref{subsec:Evaluation-metrics} for the Mexican data
across poverty lines. Specifically, these lines correspond to percentiles
that are multiples of 5 taken from the full population. }\textbf{\footnotesize{}OLS
ignoring non-normal}{\footnotesize{}: uses the method described in
Eq. \ref{eq:Fixed} to generate predictions. }\textbf{\footnotesize{}OLS
BS 20}{\footnotesize{}: Estimates obtained using }\texttt{\footnotesize{}hetmireg,
}{\footnotesize{}OLS fit across 20 bootstrap samples of training sample
with errors drawn from the empirical distribution across bootstrap
samples. }\textbf{\footnotesize{}OLS Empirical}{\footnotesize{}: uses
the method described in Eq. \ref{eq:Fixed}, but errors are drawn
from the empirical distribution of the training sample. }\textbf{\footnotesize{}MI
BS}{\footnotesize{} }\textbf{\footnotesize{}20}{\footnotesize{}: Uses
Stata's bootstrap option for multiple imputation with 20 imputations.
}\textbf{\footnotesize{}Lasso Empirical: }{\footnotesize{}applies
parameters from lasso model where $\lambda$ is selected using 10
fold cross validation, and errors are drawn from the empirical distribution
of the training sample. }\textbf{\footnotesize{}RF empirical}{\footnotesize{}:
applies parameter estimates from a random forest model to the target
data, and errors are drawn from the empirical distribution of the
training sample.}{\footnotesize\par}
\end{figure}

\begin{figure}[H]
\caption{Bias and MSE in FGT0 under non-normal errors \label{fig:sim_dgp_nonnormal_error}}

\begin{centering}
\subfloat[Bias under DGP \ref{subsec:DGP-1--} (left) and DGP \ref{subsec:DGP-2--}]{\begin{centering}
\includegraphics[scale=0.62]{C:/Users/WB378870/GitHub/S2S_RTM/5.figures/1.Normal/method_comp_fgt0_t_bias}\includegraphics[scale=0.62]{C:/Users/WB378870/GitHub/S2S_RTM/5.figures/2.Nonnormal/method_comp_fgt0_t_bias}
\par\end{centering}

}
\par\end{centering}
\begin{centering}
\subfloat[MSE under DGP \ref{subsec:DGP-1--} (left) and DGP \ref{subsec:DGP-2--}]{\begin{centering}
\includegraphics[scale=0.62]{C:/Users/WB378870/GitHub/S2S_RTM/5.figures/1.Normal/method_comp_fgt0_t_mse}\includegraphics[scale=0.62]{C:/Users/WB378870/GitHub/S2S_RTM/5.figures/2.Nonnormal/method_comp_fgt0_t_mse}
\par\end{centering}
}
\par\end{centering}
{\footnotesize{}Note: Data are generated as described in \ref{subsec:Creating-populations},
but errors are simulated from a Student's t-distribution with 10 degrees
of freedom and scaled for a SD of 0.5. Bias is assessed at various
poverty lines across the welfare distribution. Specifically, these
lines correspond to percentiles that are multiples of 5. }\textbf{\footnotesize{}Fixed
B}{\footnotesize{}: uses the method described in Eq. \ref{eq:Fixed}
to generate predictions. For each of 100 Monte Carlo (MC) simulations,
poverty is estimated at each threshold and then averaged across simulations
to produce the final estimate. }\textbf{\footnotesize{}Fixed B lnskew}{\footnotesize{}:
follows the same approach as ``Fixed B'' but the dependent variable
is transformed using a zero-skewness log transformation. }\textbf{\footnotesize{}Fixed
B bcskew}{\footnotesize{}: follows the same approach as ``Fixed B''
but the dependent variable is transformed using a Box-Cox transformation.
}\textbf{\footnotesize{}MI X}{\footnotesize{}: implements X imputations
following Eq. \ref{eq:MI}, poverty is estimated for each threshold
under each imputed vector, and the results are averaged across simulations
to obtain the final estimate. }\textbf{\footnotesize{}MI X BS}{\footnotesize{}:
same as ``MI X'' but uses Stata's bootstrap option. }\textbf{\footnotesize{}OLS
BS}{\footnotesize{}: Fits the model on bootstrap samples drawn using
simple random sampling (SRS), applies the estimated parameters to
the target data, and draws residuals from their empirical distribution.
}\textbf{\footnotesize{}lasso BIC}{\footnotesize{}: applies parameters
from lasso model where $\lambda$ is selected using the Bayesian information
criterion (BIC) function and residuals for each imputation are randomly
drawn from the predicted residuals. }\textbf{\footnotesize{}Random
Forest:}{\footnotesize{} applies parameter estimates from a random
forest model to the target data, with residuals for each imputation
randomly drawn from the distribution of predicted residuals.}{\footnotesize\par}
\end{figure}

\begin{figure}[H]
\caption{Bias and MSE in FGT0 under heteroskedasticity \label{fig:sim_dgp_het}}

\begin{centering}
\subfloat[Bias under DGP \ref{subsec:DGP-1--} (left) and DGP \ref{subsec:DGP-2--}
(right)]{\begin{centering}
\includegraphics[scale=0.62]{C:/Users/WB378870/GitHub/S2S_RTM/5.figures/1.Normal/method_comp_fgt0_het_bias}\includegraphics[scale=0.62]{C:/Users/WB378870/GitHub/S2S_RTM/5.figures/2.Nonnormal/method_comp_fgt0_het_bias}
\par\end{centering}
}
\par\end{centering}
\begin{centering}
\subfloat[MSE under DGP \ref{subsec:DGP-1--} (left) and DGP \ref{subsec:DGP-2--}]{\begin{centering}
\includegraphics[scale=0.62]{C:/Users/WB378870/GitHub/S2S_RTM/5.figures/1.Normal/method_comp_fgt0_het_mse}\includegraphics[scale=0.62]{C:/Users/WB378870/GitHub/S2S_RTM/5.figures/2.Nonnormal/method_comp_fgt0_het_mse}
\par\end{centering}
}
\par\end{centering}
{\footnotesize{}Note: Data are generated as described in \ref{subsec:Creating-populations}.
Bias is assessed at various poverty lines across the welfare distribution.
Specifically, these lines correspond to percentiles that are multiples
of 5. }\textbf{\footnotesize{}Fixed B}{\footnotesize{}: uses the method
described in Eq. \ref{eq:Fixed} to generate predictions. For each
of 100 Monte Carlo (MC) simulations, poverty is estimated at each
threshold and then averaged across simulations to produce the final
estimate. }\textbf{\footnotesize{}Fixed B lnskew}{\footnotesize{}:
follows the same approach as ``Fixed B'' but the dependent variable
is transformed using a zero-skewness log transformation. }\textbf{\footnotesize{}Fixed
B bcskew}{\footnotesize{}: follows the same approach as ``Fixed B''
but the dependent variable is transformed using a Box-Cox transformation.
}\textbf{\footnotesize{}MI X}{\footnotesize{}: implements X imputations
following Eq. \ref{eq:MI}, poverty is estimated for each threshold
under each imputed vector, and the results are averaged across simulations
to obtain the final estimate. }\textbf{\footnotesize{}MI X BS}{\footnotesize{}:
same as ``MI X'' but uses Stata's bootstrap option. }\textbf{\footnotesize{}lasso
BIC}{\footnotesize{}: applies parameters from lasso model where $\lambda$
is selected using the Bayesian information criterion (BIC) function
and residuals for each imputation are randomly drawn from the predicted
residuals. }\textbf{\footnotesize{}Het. MLE}{\footnotesize{}: uses
parameters from the \citet{harvey1976estimating} model for heteroskedasticity
and follows Eq. \ref{eq:Fixed} for the predictions where 100 MC simulations
are undertaken except that $\sigma_{e_{i}}^{2}$ is household specific.
}\textbf{\footnotesize{}Alpha model}{\footnotesize{}: obtains parameters
following \citet{elbers2002micro} model for heteroskedasticity and
follows Eq. \ref{eq:MI}, except that $\sigma_{e_{i}}^{2}$ is household
specific,with 100 imputed vectors.}{\footnotesize\par}
\end{figure}

\end{document}
